%% xiPLN: Correct, Complete PLN Foundations
%%
%% Render with:
%%   lualatex -shell-escape xiPLN.tex && bibtex xiPLN && lualatex -shell-escape xiPLN.tex && lualatex -shell-escape xiPLN.tex

\documentclass[]{article}
\usepackage{url}
\usepackage[hyperindex,breaklinks]{hyperref}
\IfFileExists{breakurl.sty}{\usepackage{breakurl}}{}
\usepackage{listings}
\lstset{basicstyle=\ttfamily\footnotesize,breaklines=true,frame=single}
\usepackage{float}
\restylefloat{table}
\usepackage{longtable}
\usepackage{graphicx}
\usepackage[font=small,labelfont=bf]{caption}
\usepackage[skip=0pt]{subcaption}
\usepackage{amsfonts}
\usepackage{amsmath}
\usepackage{amsthm}
\usepackage{unicode-math}
\IfFileExists{newunicodechar.sty}{
  \usepackage{newunicodechar}
  \newunicodechar{∘}{\ensuremath{\circ}}
}{}
\IfFontExistsTF{Stix Two Math}{\setmathfont{Stix Two Math}}{\setmathfont{STIX Math}}

% Compatibility fallback for older TeX bundles (e.g., Tectonic 2021-era
% kernel/hyperref hook stacks) where this hook target may be missing.
\makeatletter
\providecommand{\headerps@out}{}
\makeatother

\newtheorem{theorem}{Theorem}
\newtheorem{lemma}[theorem]{Lemma}
\newtheorem{definition}[theorem]{Definition}

\newcommand{\nuPLN}{\texorpdfstring{$\nu$}{nu}\textup{PLN}}
\newcommand{\xiPLN}{\texorpdfstring{$\xi$}{xi}\textup{PLN}}
\newcommand{\code}[1]{\texttt{#1}}
\newcommand{\codepath}[1]{\path{#1}}
\newcommand{\ty}[1]{\ensuremath{\mathtt{#1}}}
\newcommand{\Evidence}{\ty{Evidence}}
\newcommand{\JointEvidence}{\ty{JointEvidence}}
\newcommand{\WorldModel}{\ty{WorldModel}}
\newcommand{\hplus}{\ensuremath{\oplus}}

\title{\xiPLN: A Correct and Complete Foundation for Probabilistic Logic Networks\\
  \texttt{(DRAFT)}}
\author{Zar \and Codex CLI (OpenAI GPT-5.2)}

\begin{document}
\maketitle

\begin{abstract}
Probabilistic Logic Networks (PLN) couples logical structure with uncertain inference.
In practice, however, many PLN presentations mix distinct layers: (i) a probabilistic
semantic model, (ii) an evidence/weight plumbing model, and (iii) operational truth-value
views (strength, confidence, interval bounds).
This leads to ``semantic drift'': algebraic operators are treated as if they were the
logic itself, and link-level inference attempts to propagate correlation-sensitive
quantities without carrying the information required to do so correctly.

We present \xiPLN{}, a complete PLN foundation designed to be (a) mathematically correct
and (b) extensible to tractable sublayers.  The core move is to make the complete layer a
\emph{posterior-state calculus}: the proof object is a revisable world-model posterior
state, and all link/event truth values are obtained by querying that state.
We give a reference complete instance as a Dirichlet posterior over complete worlds
(\JointEvidence{}), prove that evidence extraction commutes with revision, and formalize a
no-go theorem showing that complete link evidence cannot, in general, be computed from local
per-link evidence alone.  We then explain how Bayesian-network-style sublayers arise as
restricted but tractable world-model classes, keeping PLN's operational rules as compiled
tactics with explicit assumptions.
\end{abstract}

\section{Motivation and Context}

PLN, as described by Goertzel, Ikl\'{e} et al.~\cite{goertzel2009pln}, represents uncertain
judgments using truth values such as
\emph{strength} (a probability-like estimate) and \emph{confidence} (a reliability-like
quantity).  Modern \emph{nuPLN} work (Nil Geisweiller, draft and v1) makes the key step
explicit: a complete grounding requires a global probabilistic semantics---a world model
with a joint distribution over possible worlds and a Bayesian update story.
We treat \nuPLN{} as the probabilistic reference layer and extend it in \xiPLN{} by (i)
making posterior states the proof objects, (ii) separating evidence algebra from probability
views, and (iii) connecting the complete layer to tractable Bayesian-network sublayers and
Heyting-valued foundations~\cite{nuPLN2026}.

Our experience formalizing PLN in Lean is that correctness failures almost always come from
mixing layers:
\begin{itemize}
\item treating operational truth-value operators (interval bounds, confidence heuristics) as if
  they were the logic; or
\item trying to compute complete posteriors for derived links from local link truth values
  without carrying the correlation information needed to do so.
\end{itemize}

\xiPLN{} is a response: it isolates a \emph{complete} core calculus and makes all fast PLN
rules live as derived/compiled approximations relative to declared model classes.

\section{Three Semantic Layers (and Why They Must Not Be Conflated)}

We distinguish three layers.

\subsection{Layer 1: World-model posterior states (complete semantics)}

The complete layer carries a posterior state in some explicit class of world models.
Conceptually: \emph{a distribution or structured distribution} that can answer any query by
conditioning/marginalization (or by exact inference within the chosen model class).

In the Lean prototype, this layer is represented by \JointEvidence{}:
\[
  \JointEvidence(n) := \texttt{Fin}(2^n) \to \mathbb{R}_{\ge 0}^{\infty}
\]
interpreted as Dirichlet pseudo-counts over the $2^n$ complete worlds for $n$ propositional
atoms.  Revision is pointwise addition of pseudo-counts.

\subsection{Layer 2: Evidence algebra (quantale / Heyting structure)}

PLN uses an evidence carrier \Evidence{} as a compositional algebra of observations.
In our formalization, \Evidence{} is the 2D count object $(n^+,n^-)$ with:
\begin{itemize}
\item parallel aggregation \hplus{} (``revision''): add counts componentwise;
\item order (information ordering): coordinatewise $\le$;
\item rich lattice structure (complete lattice; Heyting operators exist).
\end{itemize}

This layer is \emph{not} itself the complete probabilistic semantics; it is an algebraic
semantics for evidence aggregation and logical structure.
It is particularly important for FO/HO extensions, where truth values are naturally
Heyting-valued (cf. the SatisfyingSet construction in the codebase).

\subsection{Layer 3: Operational views (strength, weight, confidence, bounds)}

Strength/weight/confidence/intervals are \emph{views} extracted from evidence.
In particular, Walley-style predictive IDM intervals and Bayesian Beta-posterior
credible intervals are distinct semantics and must remain explicit selector choices
rather than being conflated into one operator~\cite{Walley1991,Walley1996BagOfMarbles,BenavoliDeCampos2009MLEDominance}.
They are indispensable operationally, but they should not be used as if they were complete
semantic state.  \xiPLN{} adopts the mantra:

\begin{quote}
Probability is what you query; evidence is what you carry.
\end{quote}

\section{\xiPLN{} Core: World-model Revision + Query}

The complete core calculus is intentionally simple: it is the calculus of revising posterior
states plus a query interface.

\subsection{World-model interface}

In Lean, we isolate a minimal interface \WorldModel{} (posterior state + query projections):

\begin{lstlisting}
inductive PLNQuery (Atom : Type*) where
  | prop : Atom -> PLNQuery Atom
  | link : Atom -> Atom -> PLNQuery Atom
  | linkCond : List Atom -> Atom -> PLNQuery Atom

class WorldModel (State : Type*) (Query : Type*) [EvidenceType State] where
  evidence : State -> Query -> Evidence
  evidence_add : forall W1 W2 q, evidence (W1 + W2) q = evidence W1 q + evidence W2 q
\end{lstlisting}

Here \code{EvidenceType} means: revision is an additive commutative monoid on \code{State}.
Standard PLN event/link queries are represented by \code{PLNQuery Atom}; all truth-value
quantities are derived by mapping extracted \Evidence{} to strengths/weights.
See \codepath{Mettapedia/Logic/PLNWorldModel.lean}.

\subsection{Reference complete instance: Dirichlet over worlds}

For finite propositional scope (atoms \code{Fin n}), we instantiate \WorldModel{} with the
Dirichlet world-table \JointEvidence{}:
\begin{itemize}
\item \code{evidence(E, prop(A))} sums world counts where an atom is \code{true} vs \code{false};
\item \code{evidence(E, link(A,B))} sums world counts where \code{A} is true and \code{B} true vs false.
\item \code{evidence(E, linkCond([A,B],C))} sums world counts where both \code{A,B} hold and \code{C} holds vs fails.
\end{itemize}

Crucially, extraction commutes with revision:
\begin{quote}
revising joint evidence and then extracting a query is equal to extracting and revising at the
query-evidence level.
\end{quote}
Formally, the Lean prototype proves:
\[
  \mathrm{evidence}(E_1+E_2, q) = \mathrm{evidence}(E_1,q) \hplus \mathrm{evidence}(E_2,q)
\]
for all queries $q$.  See \codepath{Mettapedia/Logic/PLNJointEvidence.lean}.

\subsection{Probability views}

From \Evidence{} we obtain posterior-mean probabilities (improper-prior strength):
\[
  P(A) = \frac{\#(A)}{\#(\top)},\qquad
  P(B\mid A) = \frac{\#(A\wedge B)}{\#(A)}
\]
where $\#(\cdot)$ is a world-count sum in \JointEvidence{}.
The Lean file \codepath{Mettapedia/Logic/PLNJointEvidenceProbability.lean} proves the exact ratio
forms for these views.

\section{A Sequent-Calculus View: World Models and Links}

We present \xiPLN{} as a sequent-style system with two contexts:
\begin{itemize}
\item an \emph{evidence context} $\Gamma$, a finite multiset of posterior fragments (elements of a
  world-model state type \code{State}); and
\item a \emph{side-condition context} $\Sigma$ containing structural assumptions about the chosen
  world-model class (e.g.\ a BN DAG, positivity, d-separation facts).
\end{itemize}

\subsection{Judgments}

We use two judgment forms.
\begin{enumerate}
\item \textbf{World-model judgment:} $\Sigma;\Gamma \vdash_{\mathrm{wm}} W$ meaning ``from the
  evidence pieces in $\Gamma$, we can construct the revised posterior state $W$''.
\item \textbf{Query judgment:} $\Sigma;\Gamma \vdash q \Downarrow e$ meaning ``querying the
  revised posterior state derived from $\Gamma$ yields evidence $e$ for query $q$''.
\end{enumerate}

In the Lean prototype, the world-model judgment is deterministic: if \code{State} is an
additive commutative monoid, then the constructed posterior is simply the revision sum
$W := \sum \Gamma$.  The query judgment is then obtained by extraction via \WorldModel{}:
\[
  e := \mathrm{evidence}(W,q).
\]

\subsection{Core rules}

Let \code{State} be an \code{EvidenceType} (revision $\,+\,$ with unit $0$), and let \WorldModel{}
provide $\mathrm{evidence} : \code{State} \to \code{Query} \to \Evidence$.
The complete core calculus is:
\[
\frac{}{\,\Sigma;\varnothing \vdash_{\mathrm{wm}} 0\,}\;(\mathrm{WM\text{-}Unit})
\qquad
\frac{}{\,\Sigma;\{W\} \vdash_{\mathrm{wm}} W\,}\;(\mathrm{WM\text{-}Ev})
\qquad
\frac{\Sigma;\Gamma \vdash_{\mathrm{wm}} W \quad \Sigma;\Delta \vdash_{\mathrm{wm}} W'}
     {\,\Sigma;\Gamma \uplus \Delta \vdash_{\mathrm{wm}} (W+W')\,}\;(\mathrm{WM\text{-}Rev})
\]
\[
\frac{\Sigma;\Gamma \vdash_{\mathrm{wm}} W}
     {\,\Sigma;\Gamma \vdash q \Downarrow \mathrm{evidence}(W,q)\,}\;(\mathrm{Q\text{-}Extract})
\]

The key algebraic law is commutation of extraction with revision:
\[
  \mathrm{evidence}(W_1+W_2,q) = \mathrm{evidence}(W_1,q) \hplus \mathrm{evidence}(W_2,q),
\]
which makes the following \emph{query-revision} rule admissible:
\[
\frac{\Sigma;\Gamma \vdash q \Downarrow e \quad \Sigma;\Delta \vdash q \Downarrow e'}
     {\,\Sigma;\Gamma \uplus \Delta \vdash q \Downarrow (e \hplus e')\,}\;(\mathrm{Q\text{-}Rev})
\]

\paragraph{Links are queries.}
In this view, a PLN ``link'' $A\!\Rightarrow\!B$ is not an object-level connective; it is a query
\code{link(A,B)} against the revised world model.  The calculus does not propagate links directly;
it revises world models and answers link queries.

\subsection{Derived link calculus: classic PLN rules as admissible rewrites}

Classic PLN rules (deduction/abduction/induction) become \emph{compiled tactics} or
\emph{admissible query rewrites} relative to a world-model class with explicit side conditions
$\Sigma$.  Intuitively, they are ways to answer some link queries using other link/event queries,
without pretending to compute complete link evidence in full generality.

\paragraph{Example: deduction strength admissibility (BN case).}
Let $\mathsf{strength}(W,q)$ denote the posterior-mean view extracted from evidence:
\[
  \mathsf{strength}(W,q) := \mathrm{toStrength}(\mathrm{evidence}(W,q)).
\]
For a BN world-model class, if $\Sigma$ entails the screening-off equalities needed by the general
PLN deduction theorem (conditional independence of $C$ from $A$ given $B$ and given $\neg B$, plus
positivity), then one may rewrite a hard query $\code{linkCond([A,B],C)}$ into smaller ones:
\[
  \mathsf{strength}(W,\code{linkCond([A,B],C)}) =
    \mathrm{plnDeductionStrength}\Big(
      \mathsf{strength}(W,\code{link}(A,B)),
      \mathsf{strength}(W,\code{link}(B,C)),
      \mathsf{strength}(W,\code{prop}(B)),
      \mathsf{strength}(W,\code{prop}(C))
    \Big).
\]

\paragraph{Lean status.}
\codepath{Mettapedia/Logic/PLNBayesNetFastRules.lean} proves this admissibility in the simplest
nontrivial BN instance, the chain $A\to B\to C$: the theorem
\code{chainBN\_plnDeductionStrength\_exact} shows the PLN deduction strength formula computes the
exact $P(C\mid A)$ in the chain BN under explicit positivity side conditions.
The general BN-facing rewrite layer uses the query-level screening-off form
\[
  \code{linkCond([A,B],C)} = \code{link(B,C)}
\]
under explicit semantic obligations (see \codepath{Mettapedia/Logic/PLNBNCompilation.lean}).

\paragraph{Admissible link-rule schema (strength view).}
Let $\mathsf{strength}(e)$ denote the posterior-mean view extracted from evidence $e$.
Then, for any world-model class where $\Sigma$ entails the screening-off conditions,
the following rule is admissible:
\[
\frac{\Sigma \vdash \mathsf{ScreenOff}(A,B,C) \quad
      \Sigma;\Gamma \vdash \code{link}(A,B)\Downarrow e_{AB} \quad
      \Sigma;\Gamma \vdash \code{link}(B,C)\Downarrow e_{BC} \quad
      \Sigma;\Gamma \vdash \code{prop}(B)\Downarrow e_{B} \quad
      \Sigma;\Gamma \vdash \code{prop}(C)\Downarrow e_{C}}
     {\Sigma;\Gamma \vdash \code{linkCond([A,B],C)}\Downarrow e_{AC}}
\]
with the side-condition that
\[
  \mathsf{strength}(e_{AC}) =
    \mathrm{plnDeductionStrength}(\mathsf{strength}(e_{AB}), \mathsf{strength}(e_{BC}),
                                  \mathsf{strength}(e_{B}), \mathsf{strength}(e_{C})).
\]
This is a \emph{strength-level} rewrite: the evidence $e_{AC}$ itself is still obtained by querying
the world model, unless an explicit evidence-flow law is available for the given model class.

\section{A No-Go Theorem: ``Complete'' Link Inference Cannot Be Local}

One might hope for a sequent-calculus-like system that takes local per-link evidence
(\Evidence{} for $A$, $B$, $C$, $A\!\Rightarrow\!B$, $B\!\Rightarrow\!C$) and computes complete
evidence for $A\!\Rightarrow\!C$.  \xiPLN{} asserts that this is impossible in general without
extra assumptions: correlations live in the joint state.

\begin{theorem}[No local complete deduction rule]
There is no function
\[
  f : \Evidence^5 \to \Evidence
\]
that, for all joint evidence states $E$, computes the exact link evidence for $A\!\Rightarrow\!C$
from only the local premises $\Evidence(A),\Evidence(B),\Evidence(C),\Evidence(A\!\Rightarrow\!B),
\Evidence(B\!\Rightarrow\!C)$.
\end{theorem}

The Lean proof constructs two different joint evidence states on three atoms that agree on all
these premises but disagree on the conclusion \code{linkEvidence(A,C)}.
See \codepath{Mettapedia/Logic/PLNJointEvidenceNoGo.lean}.

\paragraph{Interpretation.}
This theorem is the formal core of ``you must carry correlations''.
It does not make fast PLN useless; it tells us exactly what fast PLN \emph{cannot} claim:
completeness for arbitrary world models without structural assumptions.

\section{Tractable Sublayers: Bayesian Networks as Restricted World Models}

The complete world-table \JointEvidence{} is exponential in $n$.  To scale, we restrict the
world-model class while preserving correctness \emph{relative to that class}.

The next natural step is a Bayesian-network-style world model:
\begin{itemize}
\item a DAG structure specifying factorization;
\item local conditional tables with Dirichlet evidence (counts) per CPT row;
\item revision = add local Dirichlet counts (conjugate update);
\item query = exact inference by variable elimination / junction tree, with complexity
  controlled by treewidth.
\end{itemize}

\paragraph{Lean status.}
We have implemented the \emph{evidence plumbing} part of this sublayer:
\begin{itemize}
\item a Boolean BN CPT query type \code{CPTQuery} (node + parent configuration),
\item a CPT posterior state \code{CPTState} storing \Evidence{} per CPT entry, and
\item an additive projection from \JointEvidence{} to CPT evidence by marginalization
  (summing compatible worlds), which commutes with revision.
\end{itemize}
See \codepath{Mettapedia/Logic/PLNBayesNetWorldModel.lean}.

We have also begun the ``fast rule exactness'' bridge in the simplest nontrivial case:
\codepath{Mettapedia/Logic/PLNBayesNetFastRules.lean} sets up the chain BN $A \to B \to C$,
proves the required sink-factorization and normalization lemmas, derives the screening-off
hypotheses, and applies \code{PLNDerivation.pln\_deduction\_from\_total\_probability} to obtain an
\emph{exactness theorem}: in the chain BN, the standard PLN deduction strength formula computes the
correct conditional probability $P(C\mid A)$ (under the explicit positivity side-conditions).

Exact BN query answering (variable elimination / junction tree) is the next step: it will provide
tractable evaluation of probabilities for larger query languages, while preserving correctness
relative to the declared BN model class.

\paragraph{PLN spirit preserved.}
Classic PLN link rules (deduction/abduction/induction) become:
\begin{itemize}
\item compiled tactics that propose BN queries, or
\item conditionally-sound lemmas under explicit assumptions (e.g. conditional independence),
\end{itemize}
rather than pretending to be globally complete link-level operators.

\section{Factor Graphs: Semantic Factorization vs Operational Control}

PLN has long used factor-graph and message-passing intuition for efficient inference.
The recent MORK/MM2+ notes make this explicit by proposing a \emph{Quantale-Annotated PLN
Factor Graph} encoding in Atomspace (variables = formulas with truth values; factors = rule
instances/potentials; messages combined with $\oplus$ and $\otimes$)~\cite{Goertzel2025PLNFactorGraphMORK,Goertzel2025MOSESMORK}.
This is the right operational substrate, but \xiPLN{} separates two distinct roles:

\subsection{(A) Semantic factor graphs (world-model factorization)}

Here factors \emph{are} the model: the world-model posterior state is represented by a
factorization of a joint distribution (BN or MRF).  In this case:
\begin{itemize}
\item factors are genuine conditional / clique potentials;
\item variable elimination / junction tree gives exact query answers
  (treewidth controls complexity);
\item belief propagation is exact on trees and approximate on loopy graphs.
\end{itemize}
Thus a factor graph is a \emph{world-model representation}, and correctness is relative
to the declared model class.

\paragraph{Positive example (semantic).}
A BN with edges $A\to B\to C$ yields factors $P(A)$, $P(B\mid A)$, $P(C\mid B)$.
If the model class is fixed to this DAG, variable elimination returns the exact
$P(C\mid A)$ and the PLN deduction strength formula is admissible as a query rewrite.

\subsection{(B) Operational factor graphs (inference control)}

Here factors are \emph{rule schemas} (deduction, induction, revision, etc.) whose local
``potentials'' are truth-value update formulas.  This aligns with the PLN chainer
and MORK PathMap indexing~\cite{Goertzel2025PLNFactorGraphMORK,Goertzel2025MORKMiner}.
However, unless each factor can be justified as a true conditional or likelihood
factor in a world model, this structure should be viewed as \emph{inference control}:
it proposes which queries to ask and which rewrites to attempt, but it does not itself
constitute the semantics.

\paragraph{Negative example (operational).}
A graph whose factors are ``deduction'' or ``abduction'' update formulas on STVs is useful for
search, but it is not automatically a probabilistic model: unless each factor corresponds to a
conditional/likelihood in some WM class, message passing is a heuristic scheduler rather than
an exact semantics.

\subsection{Alignment: when the two coincide}

The two notions align when rule factors are \emph{compiled} from a world-model class:
each rule corresponds to a conditional or likelihood factor, and message passing
implements the same queries that the world-model semantics would answer.
This is exactly the intended “fast PLN rules as compiled tactics” story:
the factor graph becomes a \emph{query plan} whose correctness is discharged
by explicit side conditions ($\Sigma$) such as d-separation or Markov properties.

\paragraph{Evidence accounting and control.}
Geodesic inference control suggests a decentralized evidence ledger to prevent
double counting during message passing and chaining~\cite{Goertzel2025GeodesicInferenceControl}.
This is naturally compatible with the \WorldModel{} interface: evidence is revised
at the world-model layer, while factor-graph message passing is used to schedule
or approximate queries, with provenance tracking guarding independence assumptions.

\section{Knuth--Skilling and Heyting Foundations}

Knuth--Skilling foundations of inference~\cite{KnuthSkilling2012} provide axioms for valuation
schemes on logical lattices, extending Cox-style plausibility calculi.
In \xiPLN{}, this plays two roles:

\begin{enumerate}
\item \emph{Probability as a valuation/view.}  At the world-model layer, probabilities arise as
  valuations on the Boolean algebra of subsets of worlds.  This is the classical Kolmogorov story
  \cite{Kolmogorov1933} instantiated for finite spaces, and it underwrites the ratio theorems
  proved in \code{PLNJointEvidenceProbability}.
\item \emph{Evidence as Heyting-valued semantics.} At the evidence/algebraic layer, \Evidence{}
  forms a rich (non-Boolean) lattice.  K\&S-style valuation rules persist but Boolean equalities
  can weaken to inequalities; see \codepath{Mettapedia/Logic/EvidenceIntuitionisticProbability.lean}.
\end{enumerate}

We also bridge K\&S additive regraduations to valuation-algebra factor graphs, so VE can be run
end-to-end on KS-valued factors without assuming probability normalization; see
\codepath{Mettapedia/ProbabilityTheory/KnuthSkilling/Bridges/ValuationAlgebra.lean} and
\codepath{Mettapedia/ProbabilityTheory/BayesianNetworks/KSFactorGraph.lean}.

An important cautionary fact---the ``totality gate''---is already formalized:
because \Evidence{} has incomparable elements, it admits no faithful order-embedding into the
reals.  Thus one should not expect a single real-valued map to capture all evidence information.
See \codepath{Mettapedia/Logic/PLN_KS_Bridge.lean}.

\section{Roadmap and Deliverables}

The core Lean artifacts supporting \xiPLN{} are:
\begin{itemize}
\item \codepath{Mettapedia/Logic/PLNWorldModel.lean}: the world-model interface and derived views;
\item \codepath{Mettapedia/Logic/PLNJointEvidence.lean}: the reference complete instance and revision-commutes-with-extraction;
\item \codepath{Mettapedia/Logic/PLNBayesNetWorldModel.lean}: BN-style CPT evidence states and the additive projection from \JointEvidence{};
\item \codepath{Mettapedia/Logic/PLNBayesNetFastRules.lean}: chain BN proof that the PLN deduction strength formula is exact in $A\to B\to C$;
\item \codepath{Mettapedia/Logic/PLNJointEvidenceProbability.lean}: posterior-mean probability views (ratio theorems);
\item \codepath{Mettapedia/Logic/PLNBNCompilation.lean}: BN query compilation, \code{linkCond} queries, and Σ-guarded screening-off rewrites;
\item \codepath{Mettapedia/Logic/PLNWorldModelITV.lean}: explicit ITV semantics layer (Bayes-normal, Bayes-exact, Walley-IDM);
\item \codepath{Mettapedia/Logic/PLNWorldModelITVHypercube.lean}: interval-selector transport and typed ITV judgment plumbing;
\item \codepath{Mettapedia/Logic/PLNWMOSLFBridgeITV.lean}: typed OSLF ITV bridge and selector-specialized coherence theorems;
\item \codepath{Mettapedia/Logic/PLNXiDerivedBNRules.lean}: typed BN constructor lifts (unit/constant/dependent/native/sort-tagged);
\item \codepath{Mettapedia/Logic/PLNCanonicalAPI.lean}: canonical one-call end-to-end endpoints and final bundle APIs;
\item \codepath{Mettapedia/OSLF/Framework/QuantaleCoherence.lean}: reusable quantale transport/coherence bundle;
\item \codepath{Mettapedia/ProbabilityTheory/BayesianNetworks/VariableElimination.lean}: exact VE engine;
\item \codepath{Mettapedia/ProbabilityTheory/BayesianNetworks/ValuationAlgebra.lean}: valuation-algebra semantics + VE correctness spine;
\item \codepath{Mettapedia/ProbabilityTheory/BayesianNetworks/VEBridge.lean}: joint-measure $\leftrightarrow$ VE probability bridge;
\item \codepath{Mettapedia/ProbabilityTheory/KnuthSkilling/Bridges/ValuationAlgebra.lean}: KS → valuation-algebra bridge;
\item \codepath{Mettapedia/Logic/PLNJointEvidenceNoGo.lean}: the local-completeness no-go theorem.
\end{itemize}

Recent ITV/typed-bridge deliverables now in the Lean stack:
\begin{enumerate}
\item explicit separation of Walley predictive intervals~\cite{Walley1996BagOfMarbles}
  from Bayesian credible intervals at the semantics layer;
\item typed WM-to-OSLF bridge theorems that expose coordinate-wise threshold transport (lower/upper/credibility/width/strength);
\item one-call composed endpoints that chain selector transport, rewrite/query transport, and quantale coherence for chain/fork/collider BN constructors.
\item intensional inheritance now has a typed WM/query layer and Solomonoff-linked semantic transport:
  \codepath{Mettapedia/Logic/PLNIntensionalWorldModel.lean},
  \codepath{Mettapedia/Logic/IntensionalInheritanceSolomonoffBridge.lean},
  and linked one-call final-bundle endpoints in
  \codepath{Mettapedia/Logic/PLNCanonicalAPI.lean}.
\end{enumerate}

Next implementation milestones:
\begin{enumerate}
\item implement exact BN query answering (variable elimination first) for prop/link queries
  over the BN/factor-graph world-model layer;
\item generalize chain exactness to arbitrary BNs via d-separation/Markov properties so
  screening-off obligations are discharged structurally;
\item add a compilation layer: PLN rule patterns $\Rightarrow$ BN query plans with explicit
  $\Sigma$ side-conditions.
\item extend query language beyond atoms to formulas/events, enabling FO/HO semantics via
  SatisfyingSets while keeping probability views as world-model queries.
\item stabilize sort-tagged native-path coordinate-specialized final-bundle aliases
  (chain/fork/collider for lower/upper/credibility/width/strength);
\item close remaining proof debt in confidence-space soundness layers so all downstream
  one-call APIs rest on fully discharged theorem dependencies.
\item \textbf{Chapter 9 future work (operational closure)}: complete the operational
  large-scale inference layer once the Lean-side MeTTa runtime stack is stable enough
  for executable chainer/runtime integration; until then, Ch.\ 9 remains theorem-level
  and counterexample-first on the Lean side.
\end{enumerate}

\section{Conclusion}

\xiPLN{} formalizes the slogan ``pass distributions, not just truth values'' in a way that is
precise enough to support soundness/completeness claims and flexible enough to support tractable
sublayers.  The complete layer is not a competing ``second PLN''; it is the reference semantics
that operational PLN rules approximate, under explicit assumptions, within declared model classes.

\appendix
\section{PLN Book Coverage Map}
\label{app:pln-book-coverage}

This appendix maps the core units of the PLN book (\cite{goertzel2009pln}) to
the current Lean formalization and indicates how each unit is represented.
Coverage labels are:
\begin{itemize}
\item \textbf{Formalized}: theorem-level core is present in the main stack.
\item \textbf{Partial}: substantial core is formalized, but some book-era variants remain open.
\item \textbf{POC}: exploratory/prototype formalization, not yet integrated as a complete layer.
\end{itemize}

\begin{longtable}{p{0.09\textwidth}p{0.25\textwidth}p{0.44\textwidth}p{0.14\textwidth}}
\textbf{Book unit} & \textbf{Core content} & \textbf{Lean coverage (where/how)} & \textbf{Status} \\
\hline
Ch.\ 1 & Introduction & Conceptual framing is encoded as an explicit semantic stack split:
world-model posterior semantics vs evidence algebra vs operational truth-value views
(this paper; \codepath{Mettapedia/Logic/PLNCore.lean},
\codepath{Mettapedia/Logic/SemanticsDecisionTree.lean}). & Formalized \\
\hline
Ch.\ 2 & Knowledge representation & Query-level KR via \code{PLNQuery}, \code{WorldModel},
\code{WorldModelSigma}, and WM rewrite judgments:
\codepath{Mettapedia/Logic/PLNWorldModel.lean},
\codepath{Mettapedia/Logic/PLNWorldModelCalculus.lean},
\codepath{Mettapedia/Logic/PLNBNCompilation.lean}. & Formalized \\
\hline
Ch.\ 3 & Experiential semantics & Evidence-as-counts carrier and revision algebra:
\codepath{Mettapedia/Logic/EvidenceClass.lean},
\codepath{Mettapedia/Logic/EvidenceQuantale.lean}; complete posterior state semantics via
\codepath{Mettapedia/Logic/PLNJointEvidence.lean}. & Formalized \\
\hline
Ch.\ 4 & Indefinite truth values & ITV core plus explicit semantics split:
\codepath{Mettapedia/Logic/PLNIndefiniteTruth.lean}
(including Walley IDM predictive interval constructors),
\codepath{Mettapedia/Logic/PLNWorldModelITV.lean},
\codepath{Mettapedia/Logic/PLNWorldModelITVHypercube.lean}. & Formalized \\
\hline
Ch.\ 5 & FO extensional inference (rules/strength formulas) & Core deduction/induction/abduction
formulas and derived rules:
\codepath{Mettapedia/Logic/PLNDeduction.lean},
\codepath{Mettapedia/Logic/PLNDerivation.lean},
\codepath{Mettapedia/Logic/PLNInferenceRules.lean},
\codepath{Mettapedia/Logic/PLNMettaTruthFunctions.lean}; BN exactness anchors:
\code{chainBN\_plnDeductionStrength\_exact},
\code{forkBN\_plnDeductionStrength\_exact} in
\codepath{Mettapedia/Logic/PLNBayesNetFastRules.lean}. & Formalized \\
\hline
Ch.\ 6 & FO extensional inference with ITV & ITV semantics are transported through WM rewrites,
typed OSLF bridge, and one-call API bundles:
\codepath{Mettapedia/Logic/PLNWMOSLFBridgeITV.lean},
\codepath{Mettapedia/Logic/PLNCanonicalAPI.lean}. & Formalized \\
\hline
Ch.\ 7 & Distributional truth values & Beta/Dirichlet distributional layer:
\codepath{Mettapedia/Logic/PLNDistributional.lean},
\codepath{Mettapedia/Logic/EvidenceBeta.lean},
\codepath{Mettapedia/Logic/EvidenceDirichlet.lean}; Kyburg higher-order flattening and
decision-equivalence bridge is theorem-level and exposed at
\codepath{Mettapedia/Logic/HigherOrder/PLNKyburgReduction.lean}
with canonical aliases in \codepath{Mettapedia/Logic/PLNCanonicalAPI.lean}. & Formalized (core + Kyburg bridge) \\
\hline
Ch.\ 8 & Error magnification & Confidence-level WM transport and OSLF-threshold grounding:
\codepath{Mettapedia/Logic/PLNErrorMagnificationGrounding.lean},
canonical one-call endpoints in \codepath{Mettapedia/Logic/PLNCanonicalAPI.lean}
(including \code{ch8\_wm\_rewrite\_oslf\_threshold\_acceptance\_of\_interval});
diagnostic/characterization layers remain in
\codepath{Mettapedia/Logic/PLNBugAnalysis.lean},
\codepath{Mettapedia/Logic/Comparison/ErrorCharacterization.lean}. & Formalized (core) \\
\hline
Ch.\ 9 & Large-scale inference strategies & Two tracks:
(i) structural BN compilation/rewrite pipeline
(\codepath{Mettapedia/Logic/PLNBNCompilation.lean},
\codepath{Mettapedia/Logic/PLNXiDerivedBNRules.lean});
(ii) premise-selection/inference-control theory
(\codepath{Mettapedia/Logic/PremiseSelectionSelectorSpec.lean},
\codepath{Mettapedia/Logic/PremiseSelectionOptimality.lean},
\codepath{Mettapedia/Logic/PremiseSelectionRankingStability.lean},
\codepath{Mettapedia/Logic/PremiseSelectionCoverage.lean}).
Counterexample-first no-go/corrected-positive index:
\codepath{Mettapedia/Logic/PLNLargeScaleInferenceCounterexamples.lean}.
Operational runtimes (execution-engine implementations and performance heuristics)
are pending operational Lean-MeTTa runtime developments in this chapter map.
& Partial (counterexample-first core formalized; operational pending Lean-MeTTa runtime) \\
\hline
Ch.\ 10 & Higher-order extensional inference & HOI$\to$FOI reduction formalized:
\codepath{Mettapedia/Logic/HigherOrder/Basic.lean},
\codepath{Mettapedia/Logic/HigherOrder/HigherOrderReduction.lean},
\codepath{Mettapedia/Logic/HigherOrder.lean}. & Formalized (core) \\
\hline
Ch.\ 11 & Crisp/fuzzy quantifiers with indefinite TVs & Quantifier semantics and weakness connection:
\codepath{Mettapedia/Logic/PLNFirstOrder.lean},
\codepath{Mettapedia/Logic/PLNFirstOrder/WeaknessConnection.lean},
\codepath{Mettapedia/Logic/PLNFirstOrder/Soundness.lean}; fuzzy quantifier algebra and ITV-coordinate
bridge:
\codepath{Mettapedia/Logic/PLNFirstOrder/FuzzyQuantifierSemantics.lean},
\codepath{Mettapedia/Logic/PLNFirstOrder/FuzzyITVBridge.lean}; theorem pointers:
\code{canary\_ch11\_rule\_family\_end\_to\_end\_ext},
\code{canary\_ch11\_rule4\_not\_equivalent\_itvPath},
\code{fuzzyIntervalHolds\_strength\_of\_lower\_upper},
\code{canary\_ch11\_itv\_strength\_interval\_of\_lower\_upper}; crisp operators:
\codepath{Mettapedia/Logic/PLNCrispEvidence.lean},
\codepath{Mettapedia/Logic/PLNNoisyOr.lean}; regression target:
\codepath{Mettapedia/Logic/PLNFirstOrder/QuantifierRegression.lean},
\codepath{scripts/check\_ch11\_quantifiers.sh}. & Formalized (core + fuzzy/QFM + regressions) \\
\hline
Ch.\ 12 & Intensional inference & Information-theoretic bridge and typed WM grounding:
\codepath{Mettapedia/Logic/IntensionalInheritance.lean},
\codepath{Mettapedia/Logic/IntensionalInheritanceSolomonoffBridge.lean},
\codepath{Mettapedia/Logic/PLNIntensionalWorldModel.lean},
semantic and linked one-call endpoints in
\codepath{Mettapedia/Logic/PLNCanonicalAPI.lean}, plus executable canary/regression targets:
\codepath{Mettapedia/Logic/PLNIntensionalCanary.lean},
\codepath{Mettapedia/Logic/PLNIntensionalRegression.lean},
\codepath{scripts/check\_ch12\_intensional.sh}. & Formalized (ASSOC core) \\
\hline
Ch.\ 13 & Aspects of inference control & Formal inference-control components via selector
specification, ranking transfer/stability, aggregation-role, and executable canary/regression
targets:
\codepath{Mettapedia/Logic/PremiseSelectionSelectorSpec.lean},
\codepath{Mettapedia/Logic/PremiseSelectionOptimality.lean},
\codepath{Mettapedia/Logic/PremiseSelectionRankingStability.lean},
\codepath{Mettapedia/Logic/PremiseSelectionCoverage.lean},
\codepath{Mettapedia/Logic/PLNInferenceControlCore.lean},
\codepath{Mettapedia/Logic/PLNInferenceControlCanary.lean},
\codepath{Mettapedia/Logic/PLNInferenceControlRegression.lean},
\codepath{scripts/check\_ch13\_inference\_control.sh}. & Formalized (core + canaries) \\
\hline
Ch.\ 14 & Temporal and causal inference & Temporal operator layer and causal profile:
\codepath{Mettapedia/Logic/PLNTemporalCausalInference.lean}; event-calculus WM grounding:
\codepath{Mettapedia/Logic/PLNProbabilisticEventCalculus.lean}; temporal OSLF/GSLT and one-call
event endpoints:
\codepath{Mettapedia/Logic/PLNCanonicalAPI.lean} (event bundle/final-bundle endpoints). & Formalized (core) \\
\hline
App.\ A & Comparison of PLN and NARS rules & Consolidated theorem bundle:
\codepath{Mettapedia/Logic/PLNNARSRuleCorrespondence.lean}, with source bridges
\codepath{Mettapedia/Logic/NARSEvidenceBridge.lean} and
\codepath{Mettapedia/Logic/NARSPLNGaloisConnection.lean}. & Formalized \\
\hline
\end{longtable}

The table is intentionally conservative: rows marked \emph{Partial} indicate that the core
semantic/theorem spine is present, while some book-era variants (or full-scale operational
integration) are still being extended.

\paragraph{Coverage delta (current).}
Relative to earlier drafts, Ch.\ 12 has moved from prototype-only status to
\emph{formalized ASSOC+PAT semantic coverage}: typed WM channel grounding, canonical
\code{AssocPatSemanticModel} packaging, selector-specialized end-to-end endpoints, and a
theorem-level no-go witness (\code{no\_binary\_mixed\_policy\_for\_assocPat\_fixture}) that rules out
collapsing the concrete ASSOC+PAT mixed channel to a universal binary extensional+ASSOC law.

\subsection{Chapter 7 Rule-to-Theorem Index}
\label{app:ch7-theorem-index}

\begin{longtable}{p{0.30\textwidth}p{0.33\textwidth}p{0.29\textwidth}}
\textbf{Chapter 7 content} & \textbf{Lean theorem(s)} & \textbf{Module} \\
\hline
Context-aware strength as Beta posterior mean coordinates \newline
\textit{Tags: [PLNBook-Ch7], [Beta], [distributional]} &
\code{strengthWith\_eq\_beta\_posterior\_meanENN},
\code{evidence\_encodes\_beta\_parameters} &
\codepath{Mettapedia/Logic/HigherOrder/PLNKyburgReduction.lean} \\
\hline
Evidence aggregation as conjugate update \newline
\textit{Tags: [PLNBook-Ch7], [conjugacy], [revision]} &
\code{hplus\_is\_bayesian\_update},
\code{evidence\_aggregation\_is\_conjugate\_update} &
\codepath{Mettapedia/Logic/HigherOrder/PLNKyburgReduction.lean},
\codepath{Mettapedia/Logic/EvidenceBeta.lean} \\
\hline
Kyburg flattening and expectation-preserving reduction \newline
\textit{Tags: [PLNBook-Ch7], [Kyburg], [higher-order], [flattening]} &
\code{kyburg\_flattening},
\code{expectation\_consistency},
\code{kyburg\_no\_advantage},
\code{kyburg\_no\_advantage\_via\_monad} &
\codepath{Mettapedia/Logic/HigherOrder/PLNKyburgReduction.lean},
\codepath{Mettapedia/ProbabilityTheory/HigherOrderProbability/KyburgFlattening.lean},
\codepath{Mettapedia/ProbabilityTheory/HigherOrderProbability/GiryMonad.lean} \\
\hline
Monad associativity for flattening (typed kernel convenience) \newline
\textit{Tags: [PLNBook-Ch7], [Kyburg], [Giry], [associativity]} &
\code{flatten\_associativity},
\code{flatten\_associativity\_kernel},
\code{ch7\_flatten\_associativity},
\code{ch7\_flatten\_associativity\_kernel} &
\codepath{Mettapedia/ProbabilityTheory/HigherOrderProbability/GiryMonad.lean},
\codepath{Mettapedia/Logic/PLNCanonicalAPI.lean} \\
\hline
De Finetti singleton bridge into Kyburg flattening \newline
\textit{Tags: [PLNBook-Ch7], [DeFinetti], [mixture]} &
\code{deFinetti\_flatten\_apply\_singleton},
\code{chapter7\_distributional\_kyburg\_bridge\_available} &
\codepath{Mettapedia/Logic/HigherOrder/PLNKyburgReduction.lean} \\
\hline
Concrete worked posterior-mean fixture (uniform prior, evidence 3:1) \newline
\textit{Tags: [PLNBook-Ch7], [worked-example], [Beta]} &
\code{chapter7\_worked\_example\_strength\_uniform\_3\_1} &
\codepath{Mettapedia/Logic/HigherOrder/PLNKyburgReduction.lean} \\
\hline
\end{longtable}

\subsection{Chapter 8 Rule-to-Theorem Index}
\label{app:ch8-theorem-index}

\begin{longtable}{p{0.30\textwidth}p{0.33\textwidth}p{0.29\textwidth}}
\textbf{Chapter 8 content} & \textbf{Lean theorem(s)} & \textbf{Module} \\
\hline
View transport under WM query equivalence (confidence / strengthWith / interpret) \newline
\textit{Tags: [PLNBook-Ch8], [WM-queryEq], [view-transport]} &
\code{WMQueryEq.to\_queryStrengthWith},
\code{WMQueryEq.to\_queryStrengthWith\_threshold},
\code{WMQueryEq.to\_queryConfidence},
\code{WMQueryEq.to\_queryConfidence\_threshold},
\code{WMQueryEq.to\_queryInterpret},
\code{WMQueryEqSigma.to\_queryStrengthWith},
\code{WMQueryEqSigma.to\_queryStrengthWith\_threshold},
\code{WMQueryEqSigma.to\_queryConfidence},
\code{WMQueryEqSigma.to\_queryConfidence\_threshold},
\code{WMQueryEqSigma.to\_queryInterpret} &
\codepath{Mettapedia/Logic/PLNWorldModelCalculus.lean} \\
\hline
Weaker bridge notions (evidence preorder + view equivalence) \newline
\textit{Tags: [PLNBook-Ch8], [bridge], [evidence-order], [viewEq]} &
\code{WMQueryEq.to\_evidenceLE},
\code{WMEvidenceLE.antisymm\_to\_WMQueryEq},
\code{WMQueryEq.to\_viewEq},
\code{WMQueryEqSigma.to\_evidenceLE},
\code{WMEvidenceLESigma.antisymm\_to\_WMQueryEqSigma},
\code{WMQueryEqSigma.to\_viewEq} &
\codepath{Mettapedia/Logic/PLNWorldModelCalculus.lean} \\
\hline
Confidence-threshold atom semantics and rewrite soundness \newline
\textit{Tags: [PLNBook-Ch8], [threshold-semantics], [OSLF]} &
\code{thresholdAtomSemOfWMQConfidence},
\code{thresholdAtomSemOfWMQSigmaConfidence},
\code{wmRewriteRule\_threshold\_atom\_confidence},
\code{wmRewriteRuleSigma\_threshold\_atom\_confidence},
\code{rewrite\_then\_queryEq\_threshold\_atom\_confidence\_sigma} &
\codepath{Mettapedia/Logic/PLNErrorMagnificationGrounding.lean} \\
\hline
Threshold-bridge boundary (disjunction/diamond) and total-order recovery \newline
\textit{Tags: [PLNBook-Ch8], [boundary], [counterexample], [totality]} &
\code{threshold\_or\_counterexample},
\code{threshold\_dia\_fails},
\code{threshold\_or\_total},
\code{threshold\_dia\_total} &
\codepath{Mettapedia/Logic/OSLFEvidenceSemantics.lean} \\
\hline
Double-damping threshold-gap characterization \newline
\textit{Tags: [PLNBook-Ch8], [double-damping], [gap]} &
\code{double\_damping\_strict\_lt},
\code{threshold\_gap\_of\_double\_damping\_of\_queries} &
\codepath{Mettapedia/Logic/PLNErrorMagnificationGrounding.lean} \\
\hline
One-call WM rewrite $\to$ OSLF truth transport $\to$ threshold acceptance \newline
\textit{Tags: [PLNBook-Ch8], [end-to-end], [canonical-API]} &
\code{ch8ThresholdAccepted},
\code{ch8\_thresholdAccepted\_of\_finalBundle},
\code{ch8\_wm\_rewrite\_oslf\_threshold\_acceptance\_of\_interval},
\code{ch8\_thresholdAccepted\_bool\_selector\_fixture},
\code{end\_to\_end\_quantale\_selector\_rewrite\_query\_threshold\_finalBundle\_of\_interval} &
\codepath{Mettapedia/Logic/PLNCanonicalAPI.lean},
\codepath{Mettapedia/Logic/PLNSelectorRewriteThresholdExamples.lean} \\
\hline
WM consequence-rule layer and Kripke deontic lifts (strength inequalities) \newline
\textit{Tags: [PLNBook-Ch8], [WMConsequenceRule], [Kripke], [modal-lift]} &
\code{WMStrengthLE},
\code{WMConsequenceRule},
\code{WMConsequenceRuleOn},
\code{WMConsequenceRuleOn.ofGlobal},
\code{wm\_rexist\_reflexive\_strength\_le},
\code{wm\_dts\_ob\_pe\_strength\_le},
\code{wmRexistReflexiveConsequenceRule},
\code{wmDtsObPeConsequenceRule} &
\codepath{Mettapedia/Logic/PLNWorldModelCalculus.lean},
\codepath{Mettapedia/Logic/PLNWorldModelKripkeConsequence.lean} \\
\hline
WM shell clarifiers (state-level permissive, query-level deterministic) \newline
\textit{Tags: [PLNBook-Ch8], [WM], [determinism], [judgment-boundary]} &
\code{WMJudgment.trivial},
\code{WMJudgment.query\_deterministic},
\code{WMJudgmentCtx.query\_deterministic} &
\codepath{Mettapedia/Logic/PLNWorldModel.lean} \\
\hline
Categorical WM transport layer (institution satisfaction + hyperdoctrine Beck-Chevalley + unified endpoint) \newline
\textit{Tags: [PLNBook-Ch8], [WM], [institution], [hyperdoctrine], [Beck-Chevalley], [transport]} &
\code{satisfactionCondition\_evidence},
\code{satisfactionCondition\_strength},
\code{satisfactionCondition\_strengthWith},
\code{satisfactionCondition\_confidence},
\code{WMInstitution.sat\_condition\_strength},
\code{beckChevalley\_transport\_exists},
\code{EndpointStatement},
\code{institution\_beckChevalley\_endpoint},
\code{endpointSurface\_of\_hyperdoctrine},
\code{ch8\_institution\_beckChevalley\_endpoint\_id\_fixture},
\code{ch8\_hol\_categorical\_wrapper\_bool\_fixture},
\code{ch8\_fol\_categorical\_consequence\_singleton\_fixture} &
\codepath{Mettapedia/Logic/PLNWorldModelInstitution.lean},
\codepath{Mettapedia/Logic/PLNWorldModelHyperdoctrine.lean},
\codepath{Mettapedia/Logic/PLNWorldModelCategoricalBridge.lean},
\codepath{Mettapedia/Logic/PLNWorldModelCategoricalRegression.lean} \\
\hline
PureKernel closed-judgment $\rightarrow$ WM obligation interface (assumption-parameterized) \newline
\textit{Tags: [PLNBook-Ch8], [PureKernel], [WM], [bridge], [canary]} &
\code{PureClosedTheoryBridge},
\code{pureClosedTheoryBridge\_to\_star},
\code{WMCategoricalEndpointSurface},
\code{PureJudgmentWMInterface.profileStepStar\_sound},
\code{pureTheoryStep\_to\_wmStrengthObligation},
\code{pureTheoryStep\_to\_wmStrengthObligation\_categorical},
\code{pureTheoryStepStar\_to\_wmStrengthObligation},
\code{pureTheoryStepStar\_to\_wmStrengthObligation\_categorical},
\code{wmConsequenceRuleOn\_of\_closed\_pureTheoryStep},
\code{wmConsequenceRuleOn\_of\_closed\_pureTheoryStep\_categorical},
\code{wmConsequenceRuleOn\_of\_closed\_pureTheoryStepStar},
\code{wmConsequenceRuleOn\_of\_closed\_pureTheoryStepStar\_categorical},
\code{canary\_betaPi\_bridge\_regression\_one\_nestedLam\_wm},
\code{canary\_betaPi\_bridge\_regression\_two\_nestedLam\_wm} &
\codepath{Mettapedia/Logic/PLNWorldModelPureKernelBridge.lean} \\
\hline
GF syntactic transport $\rightarrow$ WM obligation adapters (Pure-interface aligned) + functor boundary witness \newline
\textit{Tags: [PLNBook-Ch8], [GF], [WM], [bridge], [category], [expressivity]} &
\code{gfSyntaxHom\_to\_wmStrengthObligation},
\code{gfTreePatternEq\_to\_wmStrengthObligation},
\code{wmConsequenceRuleOn\_of\_gfSyntaxHom\_viaPureInterface},
\code{wmConsequenceRuleOn\_of\_gfTreePatternEq\_viaPureInterface},
\code{wmConsequenceRuleOn\_of\_fregeStrong\_viaPureInterface},
\code{syntaxToWMFunctor\_not\_full\_of\_constant\_evidence},
\code{canary\_gfSyntaxHom\_to\_pureStyleObligation},
\code{canary\_endToEnd\_gfRule\_to\_pureStyleObligation},
\code{canary\_fregeStrong\_to\_pureStyleObligation},
\code{canary\_notFull\_constantEvidence} &
\codepath{Mettapedia/Languages/GF/GFWMObligationAdapter.lean},
\codepath{Mettapedia/Languages/GF/GFWMObligationAdapterRegression.lean} \\
\hline
HOL/FOL implication consequence-transfer wrappers into WM consequence rules \newline
\textit{Tags: [PLNBook-Ch8], [HOL], [FOL], [WMConsequenceRule], [transfer]} &
\code{wmConsequenceRule\_of\_pointwise},
\code{wmConsequenceRuleOn\_of\_pointwise},
\code{wmConsequenceRuleOn\_of\_singletonStrengthLE},
\code{externalImplication\_iff\_singletonConsequence\_of\_sound\_complete},
\code{multiset\_strength\_le\_of\_pointwise\_categorical},
\code{multiset\_strength\_le\_of\_singletonStrengthLE\_categorical},
\code{wmConsequenceRuleOn\_of\_pointwise\_categorical},
\code{wmConsequenceRuleOn\_of\_singletonStrengthLE\_categorical},
\code{pointwiseImpliesOnTheory\_of\_consequence},
\code{pointwiseImpliesOnTheory\_iff\_singletonStrengthLEOnTheory},
\code{consequence\_iff\_singletonStrengthLEOnTheory},
\code{provable\_imp\_iff\_singletonStrengthLEOnTheory},
\code{multiset\_strength\_le\_of\_consequence},
\code{wmConsequenceRuleOn\_of\_consequence},
\code{multiset\_strength\_le\_of\_provable\_imp},
\code{wmConsequenceRuleOn\_of\_provable\_imp},
\code{multiset\_strength\_le\_of\_consequence\_categorical},
\code{wmConsequenceRuleOn\_of\_consequence\_categorical},
\code{multiset\_strength\_le\_of\_provable\_imp\_categorical},
\code{wmConsequenceRuleOn\_of\_provable\_imp\_categorical} &
\codepath{Mettapedia/Logic/PLNWorldModelHOLCompleteness.lean},
\codepath{Mettapedia/Logic/PLNWorldModelFOLCompleteness.lean} \\
\hline
Singleton iff-consequence adequacy and multiset lifting \newline
\textit{Tags: [PLNBook-Ch8], [Kripke], [Neighborhood], [adequacy], [singleton], [multiset]} &
\code{singleton\_adequacy\_strength\_one},
\code{pointwiseImplies\_iff\_singletonStrengthLE},
\code{multiset\_strength\_le\_of\_singletonStrengthLE},
\code{singleton\_consequence\_schema\_parallel},
\code{multiset\_lift\_schema\_parallel} &
\codepath{Mettapedia/Logic/PLNWorldModelKripke.lean},
\codepath{Mettapedia/Logic/PLNWorldModelNeighborhood.lean} \\
\hline
Kripke proof-theoretic closure (sound/complete implication bridge) \newline
\textit{Tags: [PLNBook-Ch8], [Kripke], [Sound], [Complete], [K], [KT], [KD], [K4]} &
\code{pointwiseImpliesOn\_iff\_singletonStrengthLEOn},
\code{provable\_imp\_iff\_singletonStrengthLEOn},
\code{provable\_imp\_iff\_singletonConsequenceOn},
\code{multiset\_strength\_le\_of\_provable\_imp},
\code{multiset\_consequence\_of\_provable\_imp},
\code{provable\_imp\_iff\_singletonStrengthLEOn\_K},
\code{provable\_imp\_iff\_singletonStrengthLEOn\_KT},
\code{provable\_imp\_iff\_singletonStrengthLEOn\_KD},
\code{provable\_imp\_iff\_singletonStrengthLEOn\_K4},
\code{multiset\_strength\_le\_of\_provable\_imp\_K},
\code{multiset\_strength\_le\_of\_provable\_imp\_KT},
\code{multiset\_strength\_le\_of\_provable\_imp\_KD},
\code{multiset\_strength\_le\_of\_provable\_imp\_K4},
\code{multiset\_ob\_pe\_strength\_le\_of\_provable},
\code{multiset\_rexist\_strength\_le\_of\_provable\_KT},
\code{multiset\_ob\_pe\_strength\_le\_of\_provable\_KD} &
\codepath{Mettapedia/Logic/PLNWorldModelKripkeCompleteness.lean} \\
\hline
Kripke $\rightarrow$ neighborhood embedding transport and governance-facing mapped consequence \newline
\textit{Tags: [PLNBook-Ch8], [Kripke], [Neighborhood], [embedding], [transport], [governance]} &
\code{evidence\_mapState\_eq},
\code{queryStrength\_mapState\_eq},
\code{singletonStrengthLEOnEmbedding\_iff\_kripke},
\code{mapState\_strength\_le\_of\_provable\_imp},
\code{mapState\_governance\_ob\_pe\_strength\_le\_of\_provable} &
\codepath{Mettapedia/Logic/PLNWorldModelKripkeNeighborhoodEmbedding.lean} \\
\hline
Canonical Kripke $\rightarrow$ neighborhood representation and mapped WM consequence corollaries \newline
\textit{Tags: [PLNBook-Ch8], [Kripke], [Neighborhood], [canonical], [representation], [governance]} &
\code{satisfies\_iff\_kripke\_toNeighborhood},
\code{pointed\_satisfies\_iff},
\code{queryStrength\_mapState\_eq\_canonical},
\code{singletonStrengthLEOn\_canonical\_iff\_kripke},
\code{mapState\_strength\_le\_of\_provable\_imp\_canonical},
\code{mapState\_governance\_ob\_pe\_strength\_le\_of\_provable\_canonical},
\code{parallel\_strength\_le\_of\_provable\_imp\_canonical},
\code{parallel\_governance\_ob\_pe\_strength\_le\_of\_provable\_canonical},
\code{parallel\_strength\_le\_of\_provable\_imp\_KT\_EMT\_canonical},
\code{parallel\_strength\_le\_of\_provable\_imp\_KD\_ED\_canonical},
\code{parallel\_governance\_rexist\_strength\_le\_KT\_EMT\_canonical},
\code{parallel\_governance\_ob\_pe\_strength\_le\_KD\_ED\_canonical} &
\codepath{Mettapedia/Logic/PLNWorldModelKripkeNeighborhoodCanonical.lean} \\
\hline
Weighted/source-aware Kripke WM specialization and weight-$1$ transfer theorems \newline
\textit{Tags: [PLNBook-Ch8], [Kripke], [weighted], [provenance], [transfer]} &
\code{weightedExpansion},
\code{weightedCountP\_toWeightOne\_eq\_countP},
\code{weightedEvidence\_toWeightOne\_eq\_kripkeEvidence},
\code{queryStrength\_toWeightOne\_eq\_kripke},
\code{queryStrength\_eq\_kripkeExpansion},
\code{weighted\_strength\_le\_of\_provable\_imp},
\code{trustedGate},
\code{trustedGate\_ob\_pe\_strength\_le\_of\_provable},
\code{wmTrustedGateObPeConsequenceRule},
\code{toWeightOne\_strength\_le\_of\_provable\_imp} &
\codepath{Mettapedia/Logic/PLNWorldModelKripkeWeighted.lean} \\
\hline
Neighborhood proof-theoretic closure (sound/complete implication bridge) \newline
\textit{Tags: [PLNBook-Ch8], [Neighborhood], [Sound], [Complete], [E], [EMN], [EMT], [ED]} &
\code{pointwiseImpliesOn\_iff\_singletonStrengthLEOn},
\code{provable\_imp\_iff\_singletonStrengthLEOn},
\code{provable\_imp\_iff\_singletonConsequenceOn},
\code{multiset\_strength\_le\_of\_provable\_imp},
\code{multiset\_consequence\_of\_provable\_imp},
\code{provable\_imp\_iff\_singletonStrengthLEOn\_E},
\code{provable\_imp\_iff\_singletonStrengthLEOn\_EMN},
\code{provable\_imp\_iff\_singletonStrengthLEOn\_EMT},
\code{provable\_imp\_iff\_singletonStrengthLEOn\_ED},
\code{multiset\_strength\_le\_of\_provable\_imp\_E},
\code{multiset\_strength\_le\_of\_provable\_imp\_EMN},
\code{multiset\_strength\_le\_of\_provable\_imp\_EMT},
\code{multiset\_strength\_le\_of\_provable\_imp\_ED} &
\codepath{Mettapedia/Logic/PLNWorldModelNeighborhoodCompleteness.lean} \\
\hline
Neighborhood modal/deontic-style consequence refinement and Kripke-parallel endpoints \newline
\textit{Tags: [PLNBook-Ch8], [Neighborhood], [WMConsequenceRuleOn], [box], [diamond], [governance], [parallel], [K], [C], [Mk], [McK]} &
\code{wmConsequenceRule\_of\_provable\_imp},
\code{wmBoxConsequenceRule\_of\_provable\_imp},
\code{wmDiaConsequenceRule\_of\_provable\_imp},
\code{wm\_rexist\_reflexive\_strength\_le\_EMT},
\code{wm\_ob\_pe\_strength\_le\_ED},
\code{wmRexistReflexiveConsequenceRuleNeighborhood},
\code{wmDtsObPeConsequenceRuleNeighborhood},
\code{wm\_k\_curried\_strength\_le\_of\_provable},
\code{wm\_c\_strength\_le\_of\_provable},
\code{wm\_mck\_strength\_le\_of\_provable},
\code{wm\_mk\_strength\_le\_of\_provable},
\code{wm\_mck\_strength\_le\_of\_axiomMcK},
\code{wm\_mk\_strength\_le\_of\_axiomMk},
\code{wmMcKConsequenceRule\_of\_axiom},
\code{wmMkConsequenceRule\_of\_axiom},
\code{wm\_k\_curried\_strength\_le\_EK},
\code{wm\_k\_curried\_strength\_le\_EMK},
\code{wm\_c\_strength\_le\_EC},
\code{wm\_c\_strength\_le\_EMC},
\code{governanceModalityToNeighborhood},
\code{wm\_governance\_rexist\_strength\_le},
\code{wm\_governance\_ob\_pe\_strength\_le},
\code{governance\_translation\_endpoint\_block},
\code{kripke\_neighborhood\_rexist\_parallel},
\code{kripke\_neighborhood\_ob\_pe\_parallel} &
\codepath{Mettapedia/Logic/PLNWorldModelNeighborhoodConsequence.lean} \\
\hline
State-indexed governance translation preservation (formula-level recursive map) \newline
\textit{Tags: [PLNBook-Ch8], [WMConsequenceRuleOn], [governance], [translation], [preservation], [recursive]} &
\code{governanceBoxActionToNeighborhood},
\code{governanceDiamondActionToNeighborhood},
\code{translateGovernanceFormulaToNeighborhood},
\code{translateGovernanceFormulaToNeighborhood\_rexistFormula},
\code{translateGovernanceFormulaToNeighborhood\_obFormula},
\code{translateGovernanceFormulaToNeighborhood\_peFormula},
\code{translateGovernanceFormulaToNeighborhood\_forbiddenBox},
\code{translateGovernanceFormulaToNeighborhood\_optionalDiamond},
\code{wm\_governance\_rexist\_strength\_le\_of\_translation},
\code{wm\_governance\_ob\_pe\_strength\_le\_of\_translation},
\code{wm\_governance\_forbidden\_shape\_of\_translation},
\code{wm\_governance\_optional\_shape\_of\_translation},
\code{governance\_translation\_endpoint\_block\_of\_formula},
\code{governance\_translation\_endpoint\_block\_of\_formula\_extended} &
\codepath{Mettapedia/Logic/PLNWorldModelNeighborhoodConsequence.lean} \\
\hline
\end{longtable}

\paragraph{Chapter 8 assumption envelope.}
For threshold transport in OSLF evidence semantics, disjunction and diamond require
the explicit side conditions recorded by the theorem row above:
\code{threshold\_or\_total} requires total-order comparability of evidence values,
and \code{threshold\_dia\_total} additionally requires finite branching.
Without these conditions, use the boundary theorems
\code{threshold\_or\_counterexample} and \code{threshold\_dia\_fails}.

\paragraph{Chapter 8 naming scope note (HOL/FOL wrappers).}
Modules \code{PLNWorldModelHOLCompleteness.lean} and
\code{PLNWorldModelFOLCompleteness.lean} retain historical
\code{*Completeness} filenames, but the listed Ch.~8 endpoints are
consequence-transfer/closure wrappers. The FOL side includes a concrete
implication bridge over theory-models; the HOL side is schema-level pending
imported concrete proof-system instances.

\paragraph{Chapter 8 categorical regression target.}
Script-level target for the new WM categorical layer:
\codepath{scripts/check\_wm\_categorical.sh}.

\subsection{Chapter 9 Counterexample-First Theorem Index}
\label{app:ch9-counterexample-index}

\paragraph{Notation used in this table.}
For a selected premise set $S$ and ground-truth dependency set $D$:
``coverage'' means $\mathrm{dependencyCoverage}(D,S)=|S \cap D|$ and
``card'' means $|S|$ (the selected-set cardinality).

\begin{longtable}{p{0.30\textwidth}p{0.33\textwidth}p{0.29\textwidth}}
\textbf{Chapter 9 content} & \textbf{Lean theorem(s)} & \textbf{Module} \\
\hline
Two-term inclusion-exclusion is not identifiable for union mass \newline
\textit{Tags: [PLNBook-Ch9], [multideduction], [identifiability-no-go]} &
\code{same\_first\_two\_terms\_different\_union},
\code{ch9\_ie\_sameFirstTwo\_differentUnion} &
\codepath{Mettapedia/Logic/PLNInclusionExclusionIdentifiability.lean},
\codepath{Mettapedia/Logic/PLNLargeScaleInferenceCounterexamples.lean} \\
\hline
No universal predictor from first-two-term statistics \newline
\textit{Tags: [PLNBook-Ch9], [no-universal-estimator]} &
\code{no\_universal\_union\_predictor\_from\_first\_two\_terms},
\code{ch9\_ie\_noUniversalPredictor},
\code{no\_single\_additive\_correction\_for\_both\_models},
\code{ch9\_ie\_noSingleAdditiveCorrection} &
\codepath{Mettapedia/Logic/PLNInclusionExclusionIdentifiability.lean},
\codepath{Mettapedia/Logic/PLNLargeScaleInferenceCounterexamples.lean} \\
\hline
Multideduction residual decomposition + indexed agreement endpoint \newline
\textit{Tags: [PLNBook-Ch9], [multideduction], [positive-core], [residual]} &
\code{unionCard\_eq\_twoTerm\_plus\_residualZ},
\code{correctedEstimate\_agrees\_of\_residualAssumption},
\code{repeatedDeductionRevision\_agrees\_of\_residualAssumption},
\code{no\_universal\_constant\_correction} &
\codepath{Mettapedia/Logic/PLNMultideductionResidual.lean} \\
\hline
Indexed n-way residual decomposition + indexed agreement family \newline
\textit{Tags: [PLNBook-Ch9], [multideduction], [n-way], [positive-core]} &
\code{unionCardN\_eq\_twoTerm\_plus\_residualNZ},
\code{correctedEstimateN\_agrees\_of\_residualAssumptionN},
\code{repeatedDeductionRevisionN\_agrees\_of\_residualAssumptionN} &
\codepath{Mettapedia/Logic/PLNMultideductionResidual.lean} \\
\hline
Trail-free deterministic update can fail to stabilize (2-cycle) \newline
\textit{Tags: [PLNBook-Ch9], [protocol-dynamics], [counterexample]} &
\code{orbit\_not\_eventually\_constant},
\code{ch9\_trailFree\_notEventualConstant} &
\codepath{Mettapedia/Logic/PLNTrailFreeDynamicsCounterexample.lean},
\codepath{Mettapedia/Logic/PLNLargeScaleInferenceCounterexamples.lean} \\
\hline
Damped SP/SPN positive convergence under eventual fresh evidence \newline
\textit{Tags: [PLNBook-Ch9], [protocol-dynamics], [SP/SPN], [convergence]} &
\code{damped\_sp\_spn\_eventually\_constant\_of\_eventual\_fresh},
\code{inconsistency\_eventually\_zero\_of\_eventual\_fresh},
\code{dampedOrbit\_noFresh\_not\_eventually\_constant} &
\codepath{Mettapedia/Logic/PLNTrailFreeDampedConvergence.lean} \\
\hline
Bounded-gap fairness reset bounds for damped SP/SPN \newline
\textit{Tags: [PLNBook-Ch9], [protocol-dynamics], [SP/SPN], [fairness], [bound]} &
\code{inconsistency\_zero\_within\_bound\_of\_boundedGapFresh},
\code{inconsistency\_zero\_within\_bound\_of\_eventually\_boundedGapFresh} &
\codepath{Mettapedia/Logic/PLNTrailFreeDampedConvergence.lean} \\
\hline
Pooling axioms are insufficient to characterize additive fusion \newline
\textit{Tags: [PLNBook-Ch9], [pooling], [non-uniqueness]} &
\code{maxPoolingOperator\_ne\_fuse},
\code{ch9\_poolingAxioms\_notUnique} &
\codepath{Mettapedia/Logic/PremiseSelectionExternalBayesianity.lean},
\codepath{Mettapedia/Logic/PLNLargeScaleInferenceCounterexamples.lean} \\
\hline
Corrected uniqueness: external-Bayes + total-additivity recovers additive pool \newline
\textit{Tags: [PLNBook-Ch9], [pooling], [corrected-positive]} &
\code{poolE\_eq\_hplus\_of\_externalBayes\_totalAdd},
\code{ch9\_poolingUnique\_of\_externalBayes\_totalAdd} &
\codepath{Mettapedia/Logic/PremiseSelectionExternalBayesianity.lean},
\codepath{Mettapedia/Logic/PLNLargeScaleInferenceCounterexamples.lean} \\
\hline
Coverage/cardinality surrogate is insufficient for risk-sensitive selection \newline
\textit{Tags: [PLNBook-Ch9], [premise-selection], [surrogate-no-go]} &
\code{sameCoverageAndCard\_differentFalsePositiveRisk},
\code{no\_universal\_risk\_predictor\_from\_coverage\_and\_card},
\code{ch9\_coverageCard\_not\_sufficient\_for\_risk},
\code{ch9\_noUniversalRiskPredictor\_fromCoverageCard} &
\codepath{Mettapedia/Logic/PremiseSelectionCoverageCounterexamples.lean},
\codepath{Mettapedia/Logic/PLNLargeScaleInferenceCounterexamples.lean} \\
\hline
Finite BRGI object + policy-sensitive optimum characterization \newline
\textit{Tags: [PLNBook-Ch9], [BRGI], [positive-core], [optimization], [policy]} &
\code{mem\_feasibleSets\_iff},
\code{exists\_optimal\_feasible\_set},
\code{exists\_optimal\_feasible\_set\_with\_policy\_characterization},
\code{optimal\_value\_eq\_min\_of\_policy},
\code{ch9\_brgi\_toy\_optimal\_value} &
\codepath{Mettapedia/Logic/PremiseSelectionBRGI.lean} \\
\hline
Graph-level BRGI object + optimality-preserving refinement wrappers \newline
\textit{Tags: [PLNBook-Ch9], [BRGI], [graph-object], [refinement], [policy]} &
\code{exists\_optimal\_feasible\_set\_brg},
\code{exists\_optimal\_feasible\_set\_brg\_with\_policy\_characterization},
\code{brg\_optimal\_value\_eq\_min\_of\_policy},
\code{ch9\_brg\_graph\_toy\_optimal\_value} &
\codepath{Mettapedia/Logic/PremiseSelectionBRGI.lean} \\
\hline
Class-packaged BN d-separation discharge (non-counterexample positive path) \newline
\textit{Tags: [PLNBook-Ch9], [BN-bridge], [assumption-packaging]} &
\code{allDiscreteCPTLocalMarkov\_of},
\code{wmqueryeq\_of\_dsep\_discharge\_via\_soundness\_allLocalMarkov},
\code{chain\_screeningOff\_wmqueryeq\_of\_dsep\_allLocalMarkov},
\code{fork\_screeningOff\_wmqueryeq\_of\_dsep\_allLocalMarkov},
\code{collider\_screeningOff\_wmqueryeq\_of\_dsep\_allLocalMarkov} &
\codepath{Mettapedia/Logic/PLNBNCompilation.lean},
\codepath{Mettapedia/Logic/PLNBNLocalMarkovPackages.lean} \\
\hline
One-call selector $\to$ rewrite $\to$ threshold composition \newline
\textit{Tags: [PLNBook-Ch9], [selector], [WM-OSLF-threshold], [composed-endpoint]} &
\code{ch9\_selector\_rewrite\_threshold\_end\_to\_end\_of\_interval},
\code{bool\_selector\_rewrite\_threshold\_end\_to\_end\_fixture} &
\codepath{Mettapedia/Logic/PLNCanonicalAPI.lean},
\codepath{Mettapedia/Logic/PLNSelectorRewriteThresholdExamples.lean} \\
\hline
Chain/fork/collider single-call concrete WM$\to$OSLF$\to$threshold endpoints \newline
\textit{Tags: [PLNBook-Ch9], [BN-bridge], [single-call], [concrete]} &
\code{xi\_deduction\_threshold\_concrete\_true\_fixture},
\code{xi\_sourceRule\_threshold\_concrete\_true\_fixture},
\code{xi\_sinkRule\_threshold\_concrete\_true\_fixture} &
\codepath{Mettapedia/Logic/PLNXiDerivedBNRules.lean} \\
\hline
Independent-cause union canary (``at least one head in $n$ biased flips'') \newline
\textit{Tags: [PLNBook-Ch9], [independent-causes], [noisy-or], [executable-canary]} &
\code{noisyOrMulti\_replicate},
\code{atLeastOneHeadProb\_eq\_noisyOr\_replicate},
\code{atLeastOneHeadProb\_coin4\_p06},
\code{atLeastOneHeadProb\_coin4\_p01} &
\codepath{Mettapedia/Logic/PLNNoisyOr.lean} \\
\hline
Stochastic experiment layer: Blackwell factorization, weighted source policies, trusted-gate transfer \newline
\textit{Tags: [PLNBook-Ch9], [experiment], [Blackwell], [weighted-policy], [trusted-gate]} &
\code{expectedUtility\_eq\_of\_blackwellFactor},
\code{optimalValue\_mono\_of\_blackwellFactor},
\code{optimalValueWeighted\_mono\_of\_blackwellFactor},
\code{optimalValuePolicy\_mono\_of\_blackwellFactor},
\code{optimalValue\_trustedGate\_mono\_of\_blackwellFactor},
\code{fixture\_expectedUtility\_eq\_lifted},
\code{fixture\_optimalValue\_mono},
\code{fixture\_trustedPolicy\_optimalValue\_mono} &
\codepath{Mettapedia/Logic/PLNWorldModelExperimentStochastic.lean},
\codepath{Mettapedia/Logic/PLNWorldModelExperimentStochasticRegression.lean} \\
\hline
\end{longtable}

\paragraph{Chapter 9 MeTTa/PeTTa Usage Snippet (theorem-name anchors only).}
\begin{lstlisting}[language=]
xi_deduction_threshold_concrete_true_fixture
xi_sourceRule_threshold_concrete_true_fixture
xi_sinkRule_threshold_concrete_true_fixture
ch9_selector_rewrite_threshold_end_to_end_of_interval
bool_selector_rewrite_threshold_end_to_end_fixture
ch9_brgi_toy_optimal_value
repeatedDeductionRevision_agrees_of_residualAssumption
repeatedDeductionRevisionN_agrees_of_residualAssumptionN
damped_sp_spn_eventually_constant_of_eventual_fresh
inconsistency_zero_within_bound_of_boundedGapFresh
ch9_brg_graph_toy_optimal_value
noisyOrMulti_replicate
atLeastOneHeadProb_eq_noisyOr_replicate
expectedUtility_eq_of_blackwellFactor
optimalValue_mono_of_blackwellFactor
optimalValuePolicy_mono_of_blackwellFactor
fixture_trustedPolicy_optimalValue_mono
\end{lstlisting}

\paragraph{Chapter 9 Session Checklist (summary-template).}
\begin{itemize}
\item Counterexample-first coverage checked: inclusion-exclusion, pooling, trail-free dynamics, coverage/card risk.
\item Positive BN bridge coverage checked: chain/fork/collider d-separation to WM/OSLF/threshold endpoints.
\item Multideduction positive core checked: explicit residual decomposition + assumption-indexed agreement endpoint.
\item Multideduction n-way extension checked: indexed-family decomposition + indexed agreement family.
\item SP/SPN positive dynamics checked: eventual-fresh damping convergence + inconsistency-to-zero endpoint.
\item SP/SPN fairness bounds checked: bounded-gap reset theorem with explicit `B+1` window.
\item BRGI object coverage checked: feasible-set gate, objective bounds, and finite optimality endpoint.
\item BRGI graph refinement checked: graph-object wrappers preserve policy-sensitive optimum characterization.
\item Independent-cause canary checked: noisy-or replicate and finite concrete points.
\item Stochastic experiment layer checked: Blackwell utility monotonicity plus weighted/trusted-source policy transfer.
\item Selector$\to$rewrite$\to$threshold endpoint checked (canonical + worked fixture).
\item Note: operational runtimes are not implemented in Lean here; this section tracks theorem-level semantic/bridge coverage.
\end{itemize}

\subsection{Chapter 10 Rule-to-Theorem Index}
\label{app:ch10-theorem-index}

\begin{longtable}{p{0.30\textwidth}p{0.33\textwidth}p{0.29\textwidth}}
\textbf{Chapter 10 content} & \textbf{Lean theorem(s)} & \textbf{Module} \\
\hline
SatisfyingSet bridge: higher-order evaluation to first-order membership \newline
\textit{Tags: [PLNBook-Ch10], [HOI-FOI-bridge], [SatisfyingSet]} &
\code{member\_eq\_evaluation} &
\codepath{Mettapedia/Logic/HigherOrder/HigherOrderReduction.lean} \\
\hline
Extensional implication as subset relation \newline
\textit{Tags: [PLNBook-Ch10], [Implication], [Subset]} &
\code{extensional\_implication\_reduces\_to\_subset} &
\codepath{Mettapedia/Logic/HigherOrder/HigherOrderReduction.lean} \\
\hline
Implication-conditioned inheritance value reduction \newline
\textit{Tags: [PLNBook-Ch10], [Implication], [Inheritance], [value-level]} &
\code{inheritance\_reflects\_implication},
\code{implication\_reduces\_to\_inheritance\_value},
\code{implication\_reduces\_to\_perfect\_inheritance},
\code{chapterImplicationPred},
\code{forAll\_himp\_to\_inheritance\_ratio\_of\_bridge} &
\codepath{Mettapedia/Logic/HigherOrder/HigherOrderReduction.lean} \\
\hline
Unconditional chapter implication-value equality is false (counterexample) \newline
\textit{Tags: [PLNBook-Ch10], [Implication], [counterexample], [boundary]} &
\code{forAll\_himp\_not\_equal\_inheritance\_ratio\_unconditional} &
\codepath{Mettapedia/Logic/HigherOrder/HigherOrderReduction.lean} \\
\hline
Unconditional chapter equivalence-value equality is false (counterexample) \newline
\textit{Tags: [PLNBook-Ch10], [Equivalence], [counterexample], [boundary]} &
\code{forAll\_equivalence\_not\_equal\_similarity\_ratio\_unconditional} &
\codepath{Mettapedia/Logic/HigherOrder/HigherOrderReduction.lean} \\
\hline
Exact self-division fixed-point characterization (four p-bit corners) \newline
\textit{Tags: [PLNBook-Ch10], [value-level], [fixed-point], [corners]} &
\code{ennreal\_self\_div\_component\_fixed\_iff},
\code{self\_div\_fixed\_iff\_pbit\_corner} &
\codepath{Mettapedia/Logic/HigherOrder/HigherOrderReduction.lean} \\
\hline
Piecewise iff closures for chapter implication/equivalence ratio endpoints \newline
\textit{Tags: [PLNBook-Ch10], [Implication], [Equivalence], [iff], [piecewise]} &
\code{forAll\_himp\_ratio\_iff\_piecewise},
\code{forAll\_equivalence\_ratio\_iff\_piecewise} &
\codepath{Mettapedia/Logic/HigherOrder/HigherOrderReduction.lean} \\
\hline
Side-condition packaging and endpoint ergonomics (single bundled assumptions) \newline
\textit{Tags: [PLNBook-Ch10], [API], [side-conditions], [ergonomics]} &
\code{HimpRatioSideConditions},
\code{EquivalenceRatioSideConditions},
\code{himp\_sideConditions\_of\_true\_or\_bot},
\code{equivalence\_sideConditions\_of\_true\_or\_bot},
\code{forAll\_himp\_to\_inheritance\_ratio\_of\_sideConditions},
\code{forAll\_equivalence\_to\_similarity\_ratio\_of\_sideConditions} &
\codepath{Mettapedia/Logic/HigherOrder/HigherOrderReduction.lean} \\
\hline
Inference-rule wrappers tied to ratio-endpoint soundness \newline
\textit{Tags: [PLNBook-Ch10], [inference-rule], [soundness], [API]} &
\code{HimpReductionRule},
\code{HimpReductionRule.conclusion},
\code{HimpReductionRule.sound},
\code{EquivalenceReductionRule},
\code{EquivalenceReductionRule.conclusion},
\code{EquivalenceReductionRule.sound} &
\codepath{Mettapedia/Logic/HigherOrder/HigherOrderReduction.lean} \\
\hline
Equivalence-conditioned similarity value reduction (first proved bridge) \newline
\textit{Tags: [PLNBook-Ch10], [Equivalence], [Similarity], [value-level]} &
\code{HOEquivalent},
\code{chapterEquivalencePred},
\code{equivalence\_reduces\_to\_similarity\_value},
\code{forAll\_equivalence\_to\_similarity\_ratio\_of\_bridge},
\code{similarity\_symmetric} &
\codepath{Mettapedia/Logic/HigherOrder/HigherOrderReduction.lean} \\
\hline
Constructive corner fixtures and worked endpoint instantiation \newline
\textit{Tags: [PLNBook-Ch10], [fixtures], [corners], [worked-example]} &
\code{corner\_source\_unit\_true},
\code{corner\_source\_unit\_bot},
\code{ch10\_implication\_ratio\_fixture\_unit\_false},
\code{ch10\_equivalence\_ratio\_fixture\_unit\_both} &
\codepath{Mettapedia/Logic/HigherOrder/HigherOrderReduction.lean} \\
\hline
HOJ/HOS/Hyp/Context micro-semantics and soundness layer \newline
\textit{Tags: [PLNBook-Ch10], [HOJ], [HOS], [Hyp], [Context], [soundness]} &
\code{HigherOrderStatement},
\code{HigherOrderJudgment},
\code{hosSemantics},
\code{hojSemantics},
\code{hojAlignment},
\code{hoj\_sound\_of\_alignment},
\code{HypAtom},
\code{andExtSet},
\code{ContextMacroPayload},
\code{expandContextMacro},
\code{context\_macro\_semantics\_sound},
\code{TruthCleanPred},
\code{hosSemantics\_andExt\_of\_srcTruthClean},
\code{context\_macro\_sound\_of\_hos\_srcTruthClean},
\code{context\_macro\_sound\_of\_hoj\_alignment\_srcTruthClean} &
\codepath{Mettapedia/Logic/HigherOrder/HigherOrderReduction.lean} \\
\hline
Concrete Context/Hyp fixtures (math/juggling + mismatch boundary) \newline
\textit{Tags: [PLNBook-Ch10], [Context], [Hyp], [fixtures], [boundary]} &
\code{BenContextObj},
\code{ch10\_ctx\_payload\_math},
\code{ch10\_ctx\_payload\_juggling},
\code{ch10\_ctx\_payload\_mismatch},
\code{ch10\_context\_hyp\_fixture\_math\_sound},
\code{ch10\_context\_hyp\_fixture\_juggling\_sound},
\code{ch10\_context\_hyp\_fixture\_mismatch\_not\_sound} &
\codepath{Mettapedia/Logic/HigherOrder/HigherOrderReduction.lean} \\
\hline
Subset-conditioned inheritance structural collapse \newline
\textit{Tags: [PLNBook-Ch10], [Subset], [Inheritance], [structural]} &
\code{subset\_implies\_strong\_inheritance} &
\codepath{Mettapedia/Logic/HigherOrder/HigherOrderReduction.lean} \\
\hline
\end{longtable}

\subsection{Chapter 11 Rule-to-Theorem Index}
\label{app:ch11-theorem-index}

\begin{longtable}{p{0.30\textwidth}p{0.33\textwidth}p{0.29\textwidth}}
\textbf{Chapter 11 rule/content} & \textbf{Lean theorem(s)} & \textbf{Module} \\
\hline
Quantifier exchange (De Morgan) \newline
\textit{Tags: [PLNBook-Ch11], [Classical-DeMorgan]} &
\code{canary\_ch11\_exists\_deMorgan},
\code{canary\_ch11\_quantifier\_exchange} &
\codepath{Mettapedia/Logic/PLNFirstOrder/QuantifierCanary.lean} \\
\hline
Existential generalization (extensional) \newline
\textit{Tags: [PLNBook-Ch11], [Classical-FO]} &
\code{canary\_ch11\_existential\_generalization\_ext} &
\codepath{Mettapedia/Logic/PLNFirstOrder/QuantifierCanary.lean} \\
\hline
Universal specification (extensional) \newline
\textit{Tags: [PLNBook-Ch11], [Classical-FO]} &
\code{canary\_ch11\_universal\_specification\_ext} &
\codepath{Mettapedia/Logic/PLNFirstOrder/QuantifierCanary.lean} \\
\hline
Fuzzy existential generalization (score positivity) \newline
\textit{Tags: [PLNBook-Ch11], [Zadeh1983]} &
\code{fuzzyExistsScore\_pos\_of\_witness\_nearOne},
\code{fuzzyExistsScore\_pos\_of\_itvStrengthWitness} &
\codepath{Mettapedia/Logic/PLNFirstOrder/FuzzyQuantifierSemantics.lean},
\codepath{Mettapedia/Logic/PLNFirstOrder/FuzzyITVBridge.lean} \\
\hline
Fuzzy universal specification (`PCL = 1`) \newline
\textit{Tags: [PLNBook-Ch11], [Zadeh1983]} &
\code{nearOne\_of\_fuzzyForAll\_eq\_one},
\code{nearOne\_itvStrength\_of\_fuzzyForAll\_eq\_one} &
\codepath{Mettapedia/Logic/PLNFirstOrder/FuzzyQuantifierSemantics.lean},
\codepath{Mettapedia/Logic/PLNFirstOrder/FuzzyITVBridge.lean} \\
\hline
Rule-4 non-equivalence witness (ITV path) \newline
\textit{Tags: [PLNBook-Ch11], [ITV-counterexample]} &
\code{canary\_ch11\_rule4\_not\_equivalent\_itvPath} &
\codepath{Mettapedia/Logic/PLNFirstOrder/QuantifierWorkedExamples.lean} \\
\hline
ITV coordinate transfer (lower/upper $\Rightarrow$ strength) \newline
\textit{Tags: [PLNBook-Ch11], [Walley-IDM], [QFM-bridge]} &
\code{fuzzyIntervalHolds\_strength\_of\_lower\_upper},
\code{canary\_ch11\_itv\_strength\_interval\_of\_lower\_upper} &
\codepath{Mettapedia/Logic/PLNFirstOrder/FuzzyITVBridge.lean},
\codepath{Mettapedia/Logic/PLNFirstOrder/QuantifierWorkedExamples.lean} \\
\hline
ITV rule-family bridge (existential generalization, universal specification, exchange) \newline
\textit{Tags: [PLNBook-Ch11], [ITV-bridge], [QFM-bridge]} &
\code{fuzzyThereExistsHolds\_itvStrength\_iff\_exchange},
\code{ch11\_itv\_rule\_family\_core},
\code{canary\_ch11\_itv\_rule\_family\_bundle\_bool} &
\codepath{Mettapedia/Logic/PLNFirstOrder/FuzzyITVBridge.lean},
\codepath{Mettapedia/Logic/PLNFirstOrder/QuantifierWorkedExamples.lean} \\
\hline
Finite MOST/FEW fixtures \newline
\textit{Tags: [Zadeh1983], [Yager1984]} &
\code{canary\_ch11\_most\_fraction\_threeQuarters},
\code{canary\_ch11\_few\_fraction\_oneQuarter},
\code{canary\_ch11\_most\_vs\_few\_split\_on\_threeQuarters} &
\codepath{Mettapedia/Logic/PLNFirstOrder/QuantifierWorkedExamples.lean} \\
\hline
Zadeh-style quantified syllogism family (`MOST $\circ$ MOST`, `FEW $\circ$ MOST`) \newline
\textit{Tags: [Zadeh1983], [Yager1984]} &
\code{zadeh\_syllogism\_most\_most},
\code{zadeh\_syllogism\_few\_most},
\code{canary\_zadeh\_most\_most\_fixture},
\code{canary\_zadeh\_few\_most\_fixture} &
\codepath{Mettapedia/Logic/PLNFirstOrder/FuzzySyllogismCanary.lean} \\
\hline
QFM abstraction + generic composition transport \newline
\textit{Tags: [DiazHermida2018-QFM], [QFM-Axiomatic]} &
\code{QFMCompose},
\code{qfmMul},
\code{qfmMin},
\code{qfmLukasiewicz},
\code{qfmProbSum},
\code{qfm\_compose\_interval\_bundle},
\code{qfm\_compose\_interval\_of\_fuzzyIntervals},
\code{qfmMul\_interval\_of\_fuzzyIntervals},
\code{qfmMin\_interval\_of\_fuzzyIntervals},
\code{qfmLukasiewicz\_interval\_of\_fuzzyIntervals},
\code{qfmProbSum\_interval\_of\_fuzzyIntervals},
\code{qfm\_syllogism\_interval},
\code{qfm\_syllogism\_most\_most},
\code{qfm\_syllogism\_few\_most} &
\codepath{Mettapedia/Logic/PLNFirstOrder/FuzzyQuantifierSemantics.lean},
\codepath{Mettapedia/Logic/PLNFirstOrder/FuzzySyllogismCanary.lean} \\
\hline
QFM-oriented canaries/schemas (monotonicity, conservativity-style signature invariance) \newline
\textit{Tags: [DiazHermida2018-QFM], [Tserkovny2020]} &
\code{nearZero\_iff\_nearOne\_one\_sub},
\code{nearZeroFraction\_eq\_nearOneFraction\_one\_sub},
\code{fuzzyThereExistsHolds\_iff\_nearOneComplement},
\code{fuzzyExistsScore\_mono\_of\_pointwise},
\code{fuzzyForAllHolds\_mono\_of\_pointwise},
\code{nearOneFraction\_eq\_of\_signatureEq},
\code{fuzzyIntervalHolds\_iff\_of\_signatureEq},
\code{qfm\_composedScore\_mono\_of\_pointwise},
\code{qfm\_composedScore\_eq\_of\_signatureEq},
\code{qfm\_composedInterval\_iff\_of\_signatureEq},
\code{canary\_qfm\_monotonicity\_most},
\code{canary\_qfm\_forall\_monotonicity\_most},
\code{canary\_qfm\_conservativity\_same\_nearOne\_signature} &
\codepath{Mettapedia/Logic/PLNFirstOrder/FuzzyQuantifierSemantics.lean},
\codepath{Mettapedia/Logic/PLNFirstOrder/FuzzySyllogismCanary.lean} \\
\hline
Additional finite quantifier fixtures (`ALMOST ALL`, `ABOUT HALF`) &
\textit{Tags: [Zadeh1983], [Glöckner2006]} \newline
\code{canary\_ch11\_almostAll\_fraction\_threeQuarters},
\code{canary\_ch11\_almostAll\_holds\_on\_threeQuarters},
\code{canary\_ch11\_aboutHalf\_fraction\_oneHalf},
\code{canary\_ch11\_aboutHalf\_positive\_negative\_split} &
\codepath{Mettapedia/Logic/PLNFirstOrder/QuantifierWorkedExamples.lean} \\
\hline
Literature-calibrated verbal probability fixtures (`PROBABLY`, `MAYBE/POSSIBLY`) &
\textit{Tags: [MostellerYoutz1989], [IPCC-uncertainty-guidance]} \newline
\code{ch11ProbablyParams},
\code{canary\_ch11\_probably\_positive\_negative\_split},
\code{canary\_ch11\_probably\_holds\_on\_threeQuarters},
\code{canary\_ch11\_probably\_rejects\_oneHalf},
\code{ch11MaybeParams},
\code{canary\_ch11\_maybe\_positive\_negative\_split},
\code{canary\_ch11\_maybe\_holds\_on\_oneHalf} &
\codepath{Mettapedia/Logic/PLNFirstOrder/QuantifierWorkedExamples.lean} \\
\hline
Additional finite quantifier fixtures (`ALMOST NONE`, `MANY-BUT-NOT-MOST`) &
\textit{Tags: [Zadeh1983], [Yager1984], [DiazHermida2018-QFM]} \newline
\code{canary\_ch11\_almostNone\_fraction\_oneFifth},
\code{canary\_ch11\_almostNone\_holds\_on\_oneFifth},
\code{canary\_ch11\_manyNotMost\_fraction\_threeFifths},
\code{canary\_ch11\_manyNotMost\_vs\_most\_split} &
\codepath{Mettapedia/Logic/PLNFirstOrder/QuantifierWorkedExamples.lean} \\
\hline
QFM selector-bundle API endpoints \newline
\textit{Tags: [QFM-Axiomatic], [PLNBook-Ch11]} &
\code{qfm\_selector\_bundle\_most\_most},
\code{qfm\_selector\_bundle\_few\_most},
\code{qfm\_selector\_bundle\_most\_most\_sound},
\code{qfm\_selector\_bundle\_few\_most\_sound},
\code{qfmMin\_syllogism\_most\_most},
\code{qfmMin\_syllogism\_few\_most} &
\codepath{Mettapedia/Logic/PLNFirstOrder/FuzzySyllogismCanary.lean} \\
\hline
Non-multiplicative QFM instance theorem families \newline
\textit{Tags: [Lukasiewicz-tnorm], [Probabilistic-sum], [QFM-instance-comparison]} &
\code{qfmLukasiewicz\_syllogism\_most\_most},
\code{qfmLukasiewicz\_syllogism\_few\_most},
\code{qfmProbSum\_syllogism\_most\_most},
\code{qfmProbSum\_syllogism\_few\_most},
\code{qfm\_instance\_comparison\_most\_most},
\code{qfm\_instance\_comparison\_few\_most} &
\codepath{Mettapedia/Logic/PLNFirstOrder/FuzzySyllogismCanary.lean} \\
\hline
Regression entrypoint &
\codepath{scripts/check\_ch11\_quantifiers.sh},
\codepath{scripts/check\_ch11\_fuzzy\_syllogism.sh},
\codepath{Mettapedia/Logic/PLNFirstOrder/QuantifierRegression.lean},
\codepath{Mettapedia/Logic/PLNFirstOrder/FuzzySyllogismRegression.lean} &
\codepath{.github/workflows/lean\_action\_ci.yml} \\
\hline
\end{longtable}

\subsection{FO/HO Steelman Survey (PLN-Side)}
\label{app:fo-ho-steelman-survey}

\begin{longtable}{p{0.25\textwidth}p{0.30\textwidth}p{0.37\textwidth}}
\textbf{PLN book statement class} & \textbf{Strongest steelmanned formal form} & \textbf{Lean theorem(s)} \\
\hline
Chapter 5 FO extensional implication rules & Explicit probabilistic derivation under measurable-space assumptions (not heuristic-only formulas) & \code{pln\_deduction\_from\_total\_probability},
\code{pln\_deduction\_from\_total\_probability\_ctx},
\code{pln\_deduction\_decomposition},
\code{pln\_deduction\_uniform\_simplifies} \\
\hline
Chapter 10 HO implication/equivalence value equalities & Unconditional equalities are false; strongest valid endpoints are side-conditioned bridges plus piecewise iff closures & \code{forAll\_himp\_not\_equal\_inheritance\_ratio\_unconditional},
\code{forAll\_equivalence\_not\_equal\_similarity\_ratio\_unconditional},
\code{forAll\_himp\_to\_inheritance\_ratio\_of\_sideConditions},
\code{forAll\_equivalence\_to\_similarity\_ratio\_of\_sideConditions},
\code{forAll\_himp\_ratio\_iff\_piecewise},
\code{forAll\_equivalence\_ratio\_iff\_piecewise} \\
\hline
Chapter 10 HOJ/HOS/Hyp/Context layer & Context soundness exposed with explicit source-truth-clean side conditions and concrete positive/negative fixtures & \code{context\_macro\_sound\_of\_hoj\_alignment\_srcTruthClean},
\code{ch10\_context\_hyp\_fixture\_math\_sound},
\code{ch10\_context\_hyp\_fixture\_juggling\_sound},
\code{ch10\_context\_hyp\_fixture\_mismatch\_not\_sound} \\
\hline
Chapter 11 quantifier layer & Finite-model FO + fuzzy/QFM endpoint families with monotonicity and conservativity-style transport & \code{fuzzyExistsScore\_mono\_of\_pointwise},
\code{fuzzyForAllHolds\_mono\_of\_pointwise},
\code{qfm\_composedScore\_mono\_of\_pointwise},
\code{qfm\_composedScore\_eq\_of\_signatureEq},
\code{qfm\_composedInterval\_iff\_of\_signatureEq} \\
\hline
HOL/FOL consequence transfer into WM/PLN transport (categorical surface) & Besides concrete HOL/FOL categorical fixtures, FOL now has a full implication bridge \code{T |- (phi -> psi) <-> singleton consequence on T-models} & \code{multiset\_strength\_le\_of\_pointwise\_categorical},
\code{provable\_imp\_iff\_singletonStrengthLEOnTheory},
\code{ch8\_hol\_categorical\_wrapper\_bool\_fixture},
\code{ch8\_fol\_categorical\_consequence\_singleton\_fixture} \\
\hline
\end{longtable}

\paragraph{Scope note (Chapter 11).}
Current Chapter 11 quantifier theorems are intentionally finite-model
(\code{[Fintype U]}) results. The strongest present formal claims are finite
FO/fuzzy-QFM statements above; infinitary FO/HOL quantifier extensions remain
separate future layers.

\subsection{Chapter 12 Rule-to-Theorem Index}
\label{app:ch12-theorem-index}

\begin{longtable}{p{0.30\textwidth}p{0.33\textwidth}p{0.29\textwidth}}
\textbf{Chapter 12 content} & \textbf{Lean theorem(s)} & \textbf{Module} \\
\hline
Information-theoretic bridge to universal-mixture conditionals \newline
\textit{Tags: [PLNBook-Ch12], [Solomonoff], [log-ratio]} &
\code{intensionalFromConditional\_eq\_log2\_ratio},
\code{intensionalFromXiSemimeasure\_eq\_log2\_ratio},
\code{intensionalFromXiGeom\_eq\_log2\_ratio} &
\codepath{Mettapedia/Logic/IntensionalInheritanceSolomonoffBridge.lean} \\
\hline
Typed WM ASSOC/PAT channel grounding \newline
\textit{Tags: [PLNBook-Ch12], [WM-typed], [OSLF-bridge]} &
\code{AssocScoreCorrespondence},
\code{PATScoreCorrespondence},
\code{mixedPolicyAssoc\_of\_assocScoreCorrespondence},
\code{mixedPolicyAssocPat\_of\_assocPatScoreCorrespondence} &
\codepath{Mettapedia/Logic/PLNIntensionalWorldModel.lean} \\
\hline
Canonical semantic packaging for concrete mixed policies \newline
\textit{Tags: [PLNBook-Ch12], [ASSOC-core], [semantic-packaging]} &
\code{AssocSemanticModel},
\code{AssocPatSemanticModel},
\code{mixedPolicyAssoc\_of\_assocSemanticModel},
\code{mixedPolicyAssocPat\_of\_assocPatSemanticModel},
\code{mixedRewriteRule\_assocSemantic\_of\_model},
\code{mixedRewriteRule\_assocPatSemantic\_of\_model} &
\codepath{Mettapedia/Logic/PLNIntensionalWorldModel.lean} \\
\hline
Threshold-truth transport (semantic linked strong endpoint) \newline
\textit{Tags: [PLNBook-Ch12], [ITV], [selector-transport]} &
\code{intensional\_mixed\_assoc\_threshold\_atom\_of\_interval\_of\_solomonoff\_semantic\_linked\_strong},
\code{end\_to\_end\_quantale\_selector\_rewrite\_query\_threshold\_finalBundle\_inheritance\_solomonoff\_semantic\_linkedStrong\_of\_interval},
\code{end\_to\_end\_quantale\_selector\_rewrite\_query\_threshold\_finalBundle\_inheritance\_solomonoff\_semantic\_linkedStrong\_bayesNormal},
\code{end\_to\_end\_quantale\_selector\_rewrite\_query\_threshold\_finalBundle\_inheritance\_solomonoff\_semantic\_linkedStrong\_bayesExact},
\code{end\_to\_end\_quantale\_selector\_rewrite\_query\_threshold\_finalBundle\_inheritance\_solomonoff\_semantic\_linkedStrong\_walley} &
\codepath{Mettapedia/Logic/PLNCanonicalAPI.lean} \\
\hline
Linked-strong selector-specialized theorem wrappers \newline
\textit{Tags: [PLNBook-Ch12], [wrapper], [bayesNormal/bayesExact/walley]} &
\code{intensional\_mixed\_assoc\_threshold\_atom\_bayesNormal\_of\_solomonoff\_semantic\_linked\_strong},
\code{intensional\_mixed\_assoc\_threshold\_atom\_bayesExact\_of\_solomonoff\_semantic\_linked\_strong},
\code{intensional\_mixed\_assoc\_threshold\_atom\_walley\_of\_solomonoff\_semantic\_linked\_strong} &
\codepath{Mettapedia/Logic/PLNCanonicalAPI.lean} \\
\hline
Concrete ASSOC+PAT semantic fixture endpoints \newline
\textit{Tags: [PLNBook-Ch12], [ASSOC+PAT], [concrete-end-to-end]} &
\code{end\_to\_end\_assocPat\_bayesNormal\_lower},
\code{end\_to\_end\_assocPat\_bayesExact\_lower},
\code{end\_to\_end\_assocPat\_walley\_lower} &
\codepath{Mettapedia/Logic/PLNIntensionalAssocPatClosure.lean} \\
\hline
Canonicality no-go witness for binary collapse \newline
\textit{Tags: [PLNBook-Ch12], [ASSOC+PAT], [no-go], [canonicality]} &
\code{no\_binary\_mixed\_policy\_for\_assocPat\_fixture} &
\codepath{Mettapedia/Logic/PLNIntensionalAssocPatClosure.lean} \\
\hline
Executable canaries (positive + non-equivalence) \newline
\textit{Tags: [PLNBook-Ch12], [canary], [non-equivalence]} &
\code{canary\_ch12\_mixed\_extensional\_projection},
\code{canary\_ch12\_mixed\_assoc\_projection},
\code{canary\_ch12\_mixed\_projection\_non\_equivalent},
\code{canary\_ch12\_end\_to\_end\_bayesNormal\_lower\_threshold},
\code{canary\_ch12\_end\_to\_end\_bayesExact\_lower\_threshold},
\code{canary\_ch12\_end\_to\_end\_walley\_lower\_threshold} &
\codepath{Mettapedia/Logic/PLNIntensionalCanary.lean} \\
\hline
Regression entrypoint &
\codepath{Mettapedia/Logic/PLNIntensionalRegression.lean},
\codepath{scripts/check\_ch12\_intensional.sh} &
\codepath{Mettapedia/Logic/README.md} \\
\hline
\end{longtable}

\subsection{Chapter 13 Rule-to-Theorem Index}
\label{app:ch13-theorem-index}

\begin{longtable}{p{0.30\textwidth}p{0.33\textwidth}p{0.29\textwidth}}
\textbf{Chapter 13 content} & \textbf{Lean theorem(s)} & \textbf{Module} \\
\hline
Selector defaults + checklist linkage \newline
\textit{Tags: [PLNBook-Ch13], [selector-spec], [checklist]} &
\code{selectorSpec\_defaults\_match\_checklist},
\code{selectorSpec\_default\_priorNB\_ranking\_transfer},
\code{selectorSpec\_zeroLocalGate\_fallback} &
\codepath{Mettapedia/Logic/PremiseSelectionSelectorSpec.lean} \\
\hline
Ranking transfer from classical baselines \newline
\textit{Tags: [PLNBook-Ch13], [NB], [kNN], [ranking-transfer]} &
\code{pln\_inherits\_nb\_ranking},
\code{pln\_inherits\_knn\_ranking},
\code{pln\_knn\_ranking\_eq} &
\codepath{Mettapedia/Logic/PremiseSelectionOptimality.lean} \\
\hline
Robustness under bounded perturbations \newline
\textit{Tags: [PLNBook-Ch13], [ranking-stability], [margin]} &
\code{pairwise\_lt\_stable\_of\_margin},
\code{bayesRanking\_stable\_of\_margin\_and\_tie\_equivariant},
\code{bayesRanking\_stable\_of\_support\_bias\_margin} &
\codepath{Mettapedia/Logic/PremiseSelectionRankingStability.lean} \\
\hline
Coverage/submodularity objective family \newline
\textit{Tags: [PLNBook-Ch13], [submodularity], [greedy-coverage]} &
\code{dependencyCoverage\_mono},
\code{dependencyCoverage\_diminishing\_returns},
\code{greedyChain\_one\_minus\_exp\_bound\_sharp} &
\codepath{Mettapedia/Logic/PremiseSelectionCoverage.lean} \\
\hline
Composed Chapter-13 core endpoint (selector + stability + coverage) \newline
\textit{Tags: [PLNBook-Ch13], [composed-core], [end-to-end]} &
\code{ch13\_selector\_default\_ranking\_iff},
\code{ch13\_stable\_pooled\_of\_twoStage},
\code{ch13\_inferenceControl\_end\_to\_end} &
\codepath{Mettapedia/Logic/PLNInferenceControlCore.lean} \\
\hline
Executable greedy algorithm endpoint \newline
\textit{Tags: [PLNBook-Ch13], [algorithmic], [end-to-end]} &
\code{ch13\_inferenceControl\_end\_to\_end\_algorithmic},
\code{greedySelect\_chain\_of\_le\_card},
\code{greedySelect\_one\_minus\_exp\_bound\_of\_le\_card} &
\codepath{Mettapedia/Logic/PLNInferenceControlAlgorithms.lean} \\
\hline
Operator-style forward/backward chainer interfaces \newline
\textit{Tags: [PLNBook-Ch13], [chainer], [MeTTa/PeTTa-style]} &
\code{forwardStep},
\code{forwardSearch},
\code{backwardStep},
\code{backwardSearch},
\code{boundedSearch},
\code{forwardSearch\_chain\_of\_le\_card},
\code{boundedSearch\_eq\_forwardSearch\_of\_card\_le\_topK} &
\codepath{Mettapedia/Logic/PLNInferenceControlChainer.lean} \\
\hline
PeTTa-style propositional connection-tableau ATP model \newline
\textit{Tags: [PLNBook-Ch13], [ATP], [connection-tableau], [proof-trace-checker], [executable-DFS], [PeTTa-aligned]} &
\code{replay\_sound},
\code{replay\_complete},
\code{checkProof\_true\_iff\_root\_mem\_and\_traceDerivable},
\code{checkProof\_sound},
\code{exists\_checkProof\_true\_of\_root\_mem\_and\_connDerivable},
\code{connProveDFS\_sound},
\code{connProveDFS\_complete\_of\_trace},
\code{connProveDFS\_witness\_complete},
\code{op\_check\_proof\_sound} &
\codepath{Mettapedia/Logic/LP/PropositionalConnectionChainer.lean} \\
\hline
PeTTa fixture-parity map (\code{test\_leancop\_*} naming) \newline
\textit{Tags: [PLNBook-Ch13], [fixture-map], [parity], [ic\_unsat], [ic\_branching]} &
\code{ic\_unsat\_checkProof\_true},
\code{ic\_sat\_connProve\_none},
\code{ic\_unsat\_connProve\_has\_proof},
\code{ic\_branching\_two\_distinct\_root\_proofs},
\code{ic\_branching\_negp\_connProve\_has\_proof},
\code{ic\_branching\_negq\_connProve\_has\_proof} &
\codepath{Mettapedia/Logic/LP/PropositionalConnectionChainer.lean} \\
\hline
First-order proof-trace checker soundness (\code{proof\_fo} shape) \newline
\textit{Tags: [PLNBook-Ch13], [ATP], [FO], [proof-trace-checker], [substitution]} &
\code{traceDerivableFO\_of\_replayFO\_true},
\code{replayFO\_sound} &
\codepath{Mettapedia/Logic/LP/FirstOrderConnectionTrace.lean} \\
\hline
FO fixture map + branch-complete root-membership (PeTTa \code{fo\_*} naming) \newline
\textit{Tags: [PLNBook-Ch13], [ATP], [FO], [fixture-map], [collectProofsAll]} &
\code{fo\_unsat\_ground},
\code{fo\_unsat\_unify},
\code{fo\_sat\_singleton},
\code{proof\_mem\_collectProofsAllFO\_of\_root\_mem\_trace\_mem},
\code{fo\_unsat\_ground\_proof\_mem\_collectProofsAllFO},
\code{fo\_unsat\_unify\_proof\_mem\_collectProofsAllFO} &
\codepath{Mettapedia/Logic/LP/FirstOrderConnectionTrace.lean} \\
\hline
Operator-contract parity rows for bundle/topK outputs (\code{op\_leancop\_ch13\_bundle*}) \newline
\textit{Tags: [PLNBook-Ch13], [ATP], [operator-contract], [topK], [bundle]} &
\code{ch13Bundle\_proofs\_eq\_collectProofs},
\code{ch13Bundle\_topK\_eq\_topKCollectProofs},
\code{ch13BundleAll\_proofs\_eq\_collectProofsAll},
\code{ch13BundleAll\_topK\_eq\_topKCollectProofsAll},
\code{mem\_ch13BundleAll\_topK\_implies\_mem\_ch13BundleAll\_proofs},
\code{topKCollectProofsAllFO\_contract},
\code{mem\_topKCollectProofsAllFO\_implies\_mem\_collectProofsAllFO} &
\codepath{Mettapedia/Logic/LP/PropositionalConnectionChainer.lean},
\codepath{Mettapedia/Logic/LP/FirstOrderConnectionTrace.lean} \\
\hline
Worked finite fixtures (`topK=1` and nontrivial `topK=2`) \newline
\textit{Tags: [PLNBook-Ch13], [worked-example], [finite-fixture]} &
\code{ch13\_bool\_end\_to\_end\_algorithmic},
\code{ch13\_bool\_coverage\_exact\_one},
\code{ch13\_fin3\_end\_to\_end\_algorithmic\_topK2},
\code{ch13\_fin3\_coverage\_exact\_two} &
\codepath{Mettapedia/Logic/PLNInferenceControlExamples.lean} \\
\hline
Executable canaries (positive + negative) \newline
\textit{Tags: [PLNBook-Ch13], [canary], [non-equivalence]} &
\code{canary\_ch13\_pln\_inherits\_nb\_ranking\_bool},
\code{canary\_ch13\_pairwise\_stable\_of\_margin},
\code{canary\_ch13\_ranking\_flip\_without\_margin},
\code{canary\_ch13\_dependencyCoverage\_insert\_gain\_bool},
\code{canary\_ch13\_selector\_default\_gate\_bounds} &
\codepath{Mettapedia/Logic/PLNInferenceControlCanary.lean} \\
\hline
Regression entrypoint &
\codepath{Mettapedia/Logic/PLNInferenceControlRegression.lean},
\codepath{scripts/check\_ch13\_inference\_control.sh} &
\codepath{Mettapedia/Logic/README.md} \\
\hline
\end{longtable}

\subsection{Module-Level Coverage Index (Legacy + Recent)}
\label{app:pln-book-coverage-modules}

To make explicit that coverage is not limited to recent additions, this index lists
the principal \emph{legacy and current} modules in Mettapedia that track PLN-book content.

\paragraph{Ch.\ 3--4 (experiential semantics, indefinite truth values).}
\codepath{Mettapedia/Logic/EvidenceClass.lean},
\codepath{Mettapedia/Logic/EvidenceQuantale.lean},
\codepath{Mettapedia/Logic/PLNIndefiniteTruth.lean},
\codepath{Mettapedia/Logic/PLNWorldModelITV.lean},
\codepath{Mettapedia/Logic/PLNWorldModelITVHypercube.lean}.

\paragraph{Ch.\ 5 (FO extensional inference: rules and formulas).}
\codepath{Mettapedia/Logic/PLNDeduction.lean},
\codepath{Mettapedia/Logic/PLNDerivation.lean},
\codepath{Mettapedia/Logic/PLNInferenceRules.lean},
\codepath{Mettapedia/Logic/PLNConjunction.lean},
\codepath{Mettapedia/Logic/PLNRevision.lean},
\codepath{Mettapedia/Logic/PLNImplicantConjunction.lean},
\codepath{Mettapedia/Logic/PLNMettaTruthFunctions.lean},
\codepath{Mettapedia/Logic/PLNBayesNetFastRules.lean},
\codepath{Mettapedia/Logic/PLNBNCompilation.lean},
\codepath{Mettapedia/Logic/PLNXiDerivedBNRules.lean}.

\paragraph{Ch.\ 6--7 (ITV/distributional truth-value layers).}
\codepath{Mettapedia/Logic/PLNDistributional.lean},
\codepath{Mettapedia/Logic/EvidenceBeta.lean},
\codepath{Mettapedia/Logic/EvidenceDirichlet.lean},
\codepath{Mettapedia/Logic/HigherOrder/PLNKyburgReduction.lean},
\codepath{Mettapedia/Logic/EvidenceNormalGamma.lean},
\codepath{Mettapedia/Logic/PLNWMOSLFBridgeITV.lean},
\codepath{Mettapedia/Logic/PLNCanonicalAPI.lean}.

\paragraph{Ch.\ 8 (error magnification / diagnostics).}
\codepath{Mettapedia/Logic/PLNBugAnalysis.lean},
\codepath{Mettapedia/Logic/Comparison/ErrorCharacterization.lean},
\codepath{Mettapedia/Logic/SoundnessCompleteness.lean},
\codepath{Mettapedia/Logic/PLNSoundnessCounterexample.lean},
\codepath{Mettapedia/Logic/PLNSoundnessDiagnosis.lean}.

\paragraph{Ch.\ 9 and Ch.\ 13 (large-scale inference + inference control).}
\codepath{Mettapedia/Logic/PremiseSelectionKNN.lean},
\codepath{Mettapedia/Logic/PremiseSelectionKNN_PLNBridge.lean},
\codepath{Mettapedia/Logic/PremiseSelectionOptimality.lean},
\codepath{Mettapedia/Logic/PremiseSelectionCoverage.lean},
\codepath{Mettapedia/Logic/PremiseSelectionRankingStability.lean},
\codepath{Mettapedia/Logic/PremiseSelectionSelectorSpec.lean},
\codepath{Mettapedia/Logic/PremiseSelectionBestPLNDraft.lean},
\codepath{Mettapedia/Logic/PLNInferenceControlCore.lean},
\codepath{Mettapedia/Logic/PLNInferenceControlCanary.lean},
\codepath{Mettapedia/Logic/PLNInferenceControlRegression.lean},
\codepath{scripts/check\_ch13\_inference\_control.sh}.

\paragraph{Ch.\ 10--11 (higher-order, quantifiers, crisp/fuzzy operators).}
\codepath{Mettapedia/Logic/HigherOrder.lean},
\codepath{Mettapedia/Logic/HigherOrder/Basic.lean},
\codepath{Mettapedia/Logic/HigherOrder/HigherOrderReduction.lean},
\codepath{Mettapedia/Logic/PLNFirstOrder.lean},
\codepath{Mettapedia/Logic/PLNFirstOrder/FuzzyQuantifierSemantics.lean},
\codepath{Mettapedia/Logic/PLNFirstOrder/FuzzyITVBridge.lean},
\codepath{Mettapedia/Logic/PLNFirstOrder/QuantifierCanary.lean},
\codepath{Mettapedia/Logic/PLNFirstOrder/QuantifierWorkedExamples.lean},
\codepath{Mettapedia/Logic/PLNFirstOrder/QuantifierRegression.lean},
\codepath{Mettapedia/Logic/PLNFirstOrder/Soundness.lean},
\codepath{Mettapedia/Logic/PLNCrispEvidence.lean},
\codepath{Mettapedia/Logic/PLNNoisyOr.lean},
\codepath{Mettapedia/Logic/PLNNegation.lean},
\codepath{Mettapedia/Logic/PLNDisjunction.lean},
\codepath{scripts/check\_ch11\_quantifiers.sh}.

\paragraph{Ch.\ 12 (intensional inference).}
\codepath{Mettapedia/Logic/IntensionalInheritance.lean},
\codepath{Mettapedia/Logic/IntensionalInheritanceSolomonoffBridge.lean},
\codepath{Mettapedia/Logic/PLNIntensionalWorldModel.lean},
\codepath{Mettapedia/Logic/PLNCanonicalAPI.lean},
\codepath{Mettapedia/Logic/PLNIntensionalCanary.lean},
\codepath{Mettapedia/Logic/PLNIntensionalRegression.lean},
\codepath{scripts/check\_ch12\_intensional.sh}.

\paragraph{Ch.\ 14 (temporal/causal; event calculus; OSLF/GSLT grounding).}
\codepath{Mettapedia/Logic/PLNTemporal.lean},
\codepath{Mettapedia/Logic/TemporalQuantale.lean},
\codepath{Mettapedia/Logic/PLNTemporalCausalInference.lean},
\codepath{Mettapedia/Logic/PLNProbabilisticEventCalculus.lean},
\codepath{Mettapedia/OSLF/Framework/VertexTemporalRewriteRules.lean},
\codepath{Mettapedia/OSLF/Framework/HypercubeTemporalGSLTFunctor.lean},
\codepath{Mettapedia/OSLF/Framework/QuantaleCoherence.lean}.

\paragraph{Appendix A (PLN/NARS comparison).}
\codepath{Mettapedia/Logic/PLNNARSRuleCorrespondence.lean},
\codepath{Mettapedia/Logic/NARSEvidenceBridge.lean},
\codepath{Mettapedia/Logic/NARSPLNGaloisConnection.lean}.

This module-level index is intended as the explicit ``not just the new files'' map.

\bibliographystyle{plain}
\bibliography{references}

\end{document}
